\documentclass[11pt,a4paper]{article}
\usepackage[utf8]{inputenc}
\usepackage{lmodern}
% Latin Modern: a scalable Type 1 font. With T1 and the default
% Computer Modern, TeX asks for the EC fonts, and when the Type 1 EC
% ones (cm-super) are missing it falls back to bitmaps -- and
% microtype's font expansion cannot work with a bitmap font, so the
% build fails on a TeX Live without cm-super.
\usepackage[T1]{fontenc}
\usepackage[english]{babel}
\usepackage[margin=2.3cm]{geometry}
\usepackage{xcolor}
\usepackage{booktabs}
\usepackage{longtable}
\usepackage{array}
\usepackage{amsmath,amssymb}
\usepackage{listings}
\usepackage[hidelinks]{hyperref}
\newcommand{\code}[1]{\texttt{#1}}
% Generated by tools/gen_concepts_values.py -- do not edit.
% Values come from mal-core/config/model_constants.yaml via the
% validated loader; a key that no longer exists makes \mskey fail
% the build with an error, in any document that references it.

\makeatletter
\newcommand{\mskey}[1]{%
  \expandafter\ifx\csname ms@k@#1\endcsname\relax
    \errmessage{Concept document: unknown model constant key `#1'}%
  \else
    \csname ms@k@#1\endcsname%
  \fi}
\makeatother

\expandafter\def\csname ms@k@adults.daily_mort_basal\endcsname{0.87}
\expandafter\def\csname ms@k@adults.mort_cap\endcsname{0.93}
\expandafter\def\csname ms@k@adults.mort_floor\endcsname{0.6}
\expandafter\def\csname ms@k@adults.opt_c\endcsname{25.0}
\expandafter\def\csname ms@k@adults.sigma_c\endcsname{15.0}
\expandafter\def\csname ms@k@aquatic_stages.egg_daily_mort\endcsname{0.08}
\expandafter\def\csname ms@k@aquatic_stages.larva_bh_alpha\endcsname{0.05}
\expandafter\def\csname ms@k@aquatic_stages.larva_bh_s0\endcsname{0.95}
\expandafter\def\csname ms@k@aquatic_stages.last_larval_instar\endcsname{4}
\expandafter\def\csname ms@k@aquatic_stages.pupa_daily_mort\endcsname{0.05}
\expandafter\def\csname ms@k@aquatic_stages.salinity_missing_psu\endcsname{0.0}
\expandafter\def\csname ms@k@aquatic_stress.desicc_base_daily\endcsname{0.1}
\expandafter\def\csname ms@k@aquatic_stress.desicc_grace_days\endcsname{5.0}
\expandafter\def\csname ms@k@aquatic_stress.desicc_ramp_decay\endcsname{0.3}
\expandafter\def\csname ms@k@aquatic_stress.vulnerable_instar_max\endcsname{2}
\expandafter\def\csname ms@k@aquatic_stress.wash_fraction_max\endcsname{0.6}
\expandafter\def\csname ms@k@aquatic_stress.wash_threshold_mm\endcsname{40.0}
\expandafter\def\csname ms@k@capacity.persistent_availability\endcsname{1.0}
\expandafter\def\csname ms@k@contracts.cogs_state\endcsname{"4.0"}
\expandafter\def\csname ms@k@contracts.env_nc\endcsname{"3.0"}
\expandafter\def\csname ms@k@contracts.habitat\endcsname{"4.0"}
\expandafter\def\csname ms@k@dispersal.flight_hour\endcsname{20}
\expandafter\def\csname ms@k@dispersal.night_end_hour\endcsname{6}
\expandafter\def\csname ms@k@dispersal.night_start_hour\endcsname{18}
\expandafter\def\csname ms@k@dispersal.snap_radius_m\endcsname{500.0}
\expandafter\def\csname ms@k@dispersal.wind_flight_hours\endcsname{4.0}
\expandafter\def\csname ms@k@dispersal.wind_migration_prob\endcsname{0.05}
\expandafter\def\csname ms@k@dispersal.wind_self_speed_ms\endcsname{1.0}
\expandafter\def\csname ms@k@dispersal.wind_survival\endcsname{0.85}
\expandafter\def\csname ms@k@eip.base_c\endcsname{16.0}
\expandafter\def\csname ms@k@eip.threshold_gd\endcsname{110.0}
\expandafter\def\csname ms@k@geodesy.earth_radius_m\endcsname{6371008.8}
\expandafter\def\csname ms@k@geodesy.metres_per_deg_lat\endcsname{111195.08023353292}
\expandafter\def\csname ms@k@gonotrophic.cycle_duration_days\endcsname{2.0}
\expandafter\def\csname ms@k@gonotrophic.egg_batch_max\endcsname{170}
\expandafter\def\csname ms@k@gonotrophic.egg_batch_mean\endcsname{120}
\expandafter\def\csname ms@k@gonotrophic.egg_batch_min\endcsname{30}
\expandafter\def\csname ms@k@gonotrophic.egg_batch_n\endcsname{240}
\expandafter\def\csname ms@k@gonotrophic.feeding_success_rate\endcsname{0.825}
\expandafter\def\csname ms@k@gonotrophic.first_cycle_days\endcsname{4.0}
\expandafter\def\csname ms@k@gonotrophic.resting_duration_days\endcsname{0.5}
\expandafter\def\csname ms@k@habitat.activation_min_engine\endcsname{0.01}
\expandafter\def\csname ms@k@habitat.candidate_min_engine\endcsname{0.01}
\expandafter\def\csname ms@k@habitat_types.dynamics.crop_pluvial\endcsname{"pluvial"}
\expandafter\def\csname ms@k@habitat_types.dynamics.open_pluvial\endcsname{"pluvial"}
\expandafter\def\csname ms@k@habitat_types.dynamics.surface_persistent\endcsname{"persistent"}
\expandafter\def\csname ms@k@habitat_types.dynamics.urban_persistent\endcsname{"persistent"}
\expandafter\def\csname ms@k@habitat_types.dynamics.wetland_persistent\endcsname{"persistent"}
\expandafter\def\csname ms@k@habitat_types.dynamics.wetland_pluvial\endcsname{"pluvial"}
\expandafter\def\csname ms@k@habitat_types.types\endcsname{"surface\_persistent, wetland\_persistent, urban\_persistent, wetland\_pluvial, open\_pluvial, crop\_pluvial"}
\expandafter\def\csname ms@k@host_seeking.indoor_attraction_add\endcsname{0.72}
\expandafter\def\csname ms@k@host_seeking.indoor_fraction_rural\endcsname{0.3}
\expandafter\def\csname ms@k@host_seeking.indoor_fraction_urban\endcsname{0.72}
\expandafter\def\csname ms@k@host_seeking.starvation_days\endcsname{7}
\expandafter\def\csname ms@k@host_seeking.urbanicity_add\endcsname{0.2}
\expandafter\def\csname ms@k@ingest.candidate_min_ingest\endcsname{1e-06}
\expandafter\def\csname ms@k@ingest.cell_size_hint_m\endcsname{1000.0}
\expandafter\def\csname ms@k@ingest.coastline_buffer_m\endcsname{2000.0}
\expandafter\def\csname ms@k@ingest.coastline_resolution\endcsname{"f"}
\expandafter\def\csname ms@k@ingest.fill_frac_valid_max\endcsname{0.5}
\expandafter\def\csname ms@k@ingest.grid_degrees_step_max\endcsname{1}
\expandafter\def\csname ms@k@ingest.land_mask_threshold\endcsname{0.5}
\expandafter\def\csname ms@k@ingest.min_dilation_radius_cells\endcsname{1}
\expandafter\def\csname ms@k@ingest.nc_compression_level\endcsname{4}
\expandafter\def\csname ms@k@ingest.nodata_sentinel\endcsname{-9999.0}
\expandafter\def\csname ms@k@ingest.nodata_sentinel_int\endcsname{-9999}
\expandafter\def\csname ms@k@ingest.ocean_mask_threshold\endcsname{0.5}
\expandafter\def\csname ms@k@ingest.residential_from_indoor_frac\endcsname{0.8}
\expandafter\def\csname ms@k@ingest.sea_water_floor\endcsname{0.05}
\expandafter\def\csname ms@k@ingest.valid_count_epsilon\endcsname{1e-09}
\expandafter\def\csname ms@k@integration.directed_flight_budget_h\endcsname{0.4856}
\expandafter\def\csname ms@k@integration.dt_olfactory_s\endcsname{0.25}
\expandafter\def\csname ms@k@integration.dt_terminal_s\endcsname{0.04}
\expandafter\def\csname ms@k@integration.landing_speed_hold_s\endcsname{3.0}
\expandafter\def\csname ms@k@integration.landing_speed_mm_s\endcsname{50.0}
\expandafter\def\csname ms@k@integration.phase_duration_hours\endcsname{"12.0, 4.0, 6.0, 2.0"}
\expandafter\def\csname ms@k@integration.plume_memory_s\endcsname{0.7}
\expandafter\def\csname ms@k@mating.female_attendance_rate\endcsname{0.1033}
\expandafter\def\csname ms@k@mating.local_swarm_male_offer_factor\endcsname{0.0}
\expandafter\def\csname ms@k@mating.male_females_per_night_max\endcsname{3}
\expandafter\def\csname ms@k@mating.male_move_capacity_m\endcsname{50.0}
\expandafter\def\csname ms@k@mating.male_swarm_reference_size\endcsname{100.0}
\expandafter\def\csname ms@k@mating.p_mate_male_fraction\endcsname{0.072}
\expandafter\def\csname ms@k@mating.p_mate_pairs_per_individual\endcsname{0.167}
\expandafter\def\csname ms@k@mating.p_mate_size_exponent\endcsname{0.0}
\expandafter\def\csname ms@k@mating.rain_suppresses_swarming\endcsname{true}
\expandafter\def\csname ms@k@mating.swarm_duration_min\endcsname{28.3}
\expandafter\def\csname ms@k@mating.swarm_height_m\endcsname{2.87}
\expandafter\def\csname ms@k@mating.swarm_markers_per_km2\endcsname{1.0}
\expandafter\def\csname ms@k@mating.swarm_markers_source\endcsname{"engineering prior"}
\expandafter\def\csname ms@k@mating.swarm_start_min_after_sunset\endcsname{5.0}
\expandafter\def\csname ms@k@oviposition.detect_range_m\endcsname{2000.0}
\expandafter\def\csname ms@k@oviposition.min_attraction\endcsname{0.3}
\expandafter\def\csname ms@k@oviposition.scale_m\endcsname{50.0}
\expandafter\def\csname ms@k@oviposition.step_m\endcsname{50.0}
\expandafter\def\csname ms@k@oviposition.timeout_days\endcsname{7}
\expandafter\def\csname ms@k@pluvial.cold_start_z\endcsname{0.0}
\expandafter\def\csname ms@k@pluvial.delta_days\endcsname{1.0}
\expandafter\def\csname ms@k@pluvial.dry_memory_days\endcsname{5.0}
\expandafter\def\csname ms@k@pluvial.rain_scale_mm\endcsname{2.5}
\expandafter\def\csname ms@k@rng.aquatic_seed_larva_mortality\endcsname{123}
\expandafter\def\csname ms@k@rng.aquatic_seed_stage_mortality\endcsname{42}
\expandafter\def\csname ms@k@rng.aquatic_seed_washout\endcsname{789}
\expandafter\def\csname ms@k@species.arabiensis.activity_weights\endcsname{"0.02, 0.1, 0.8, 0.08"}
\expandafter\def\csname ms@k@species.arabiensis.aestivation_supported\endcsname{true}
\expandafter\def\csname ms@k@species.arabiensis.anemotaxis_deviation_deg\endcsname{30.0}
\expandafter\def\csname ms@k@species.arabiensis.anemotaxis_rule\endcsname{"upwind"}
\expandafter\def\csname ms@k@species.arabiensis.arrested_daily_survival\endcsname{0.986}
\expandafter\def\csname ms@k@species.arabiensis.blood_meal_refractory_days\endcsname{0.5}
\expandafter\def\csname ms@k@species.arabiensis.blood_meal_size_ul\endcsname{0.3}
\expandafter\def\csname ms@k@species.arabiensis.blood_meal_target_ul\endcsname{0.5}
\expandafter\def\csname ms@k@species.arabiensis.casting_angle_deg\endcsname{90.0}
\expandafter\def\csname ms@k@species.arabiensis.casting_probability\endcsname{0.1}
\expandafter\def\csname ms@k@species.arabiensis.co2_activation_radius_m\endcsname{100.0}
\expandafter\def\csname ms@k@species.arabiensis.daily_move_capacity_m\endcsname{590.0}
\expandafter\def\csname ms@k@species.arabiensis.directed_flight_budget_h\endcsname{0.4856}
\expandafter\def\csname ms@k@species.arabiensis.eggs_per_ul\endcsname{223.0}
\expandafter\def\csname ms@k@species.arabiensis.exploration_range_m\endcsname{100.0}
\expandafter\def\csname ms@k@species.arabiensis.exploration_step_correlation\endcsname{0.5}
\expandafter\def\csname ms@k@species.arabiensis.exploration_step_m\endcsname{50.0}
\expandafter\def\csname ms@k@species.arabiensis.first_cycle_target_ul\endcsname{0.8}
\expandafter\def\csname ms@k@species.arabiensis.flight_speed_ms\endcsname{0.3375}
\expandafter\def\csname ms@k@species.arabiensis.gonotrophic_arrest_prob_dry\endcsname{0.4454}
\expandafter\def\csname ms@k@species.arabiensis.male_daily_survival\endcsname{0.773}
\expandafter\def\csname ms@k@species.arabiensis.male_ratio\endcsname{0.5}
\expandafter\def\csname ms@k@species.arabiensis.meal_size_min_ul\endcsname{0.1}
\expandafter\def\csname ms@k@species.arabiensis.olfactory_radius_m\endcsname{50.0}
\expandafter\def\csname ms@k@species.arabiensis.plume_memory_s\endcsname{0.7}
\expandafter\def\csname ms@k@species.arabiensis.pref.cattle\endcsname{0.2}
\expandafter\def\csname ms@k@species.arabiensis.pref.goat\endcsname{0.15}
\expandafter\def\csname ms@k@species.arabiensis.pref.human\endcsname{0.5}
\expandafter\def\csname ms@k@species.arabiensis.pref.sheep\endcsname{0.1}
\expandafter\def\csname ms@k@species.arabiensis.pref.wildlife\endcsname{0.05}
\expandafter\def\csname ms@k@species.arabiensis.probing_success_rate\endcsname{0.7275}
\expandafter\def\csname ms@k@species.arabiensis.salinity_hi_tol_ppt\endcsname{15.0}
\expandafter\def\csname ms@k@species.arabiensis.salinity_opt_ppt\endcsname{0.0}
\expandafter\def\csname ms@k@species.arabiensis.salinity_slope\endcsname{0.15}
\expandafter\def\csname ms@k@species.arabiensis.terminal_radius_m\endcsname{0.15}
\expandafter\def\csname ms@k@species.arabiensis.visual_radius_m\endcsname{0.485}
\expandafter\def\csname ms@k@species.arabiensis.weights.crop_pluvial\endcsname{0.1}
\expandafter\def\csname ms@k@species.arabiensis.weights.open_pluvial\endcsname{0.55}
\expandafter\def\csname ms@k@species.arabiensis.weights.surface_persistent\endcsname{0.1}
\expandafter\def\csname ms@k@species.arabiensis.weights.urban_persistent\endcsname{0.15}
\expandafter\def\csname ms@k@species.arabiensis.weights.wetland_persistent\endcsname{0.1}
\expandafter\def\csname ms@k@species.arabiensis.weights.wetland_pluvial\endcsname{0.0}
\expandafter\def\csname ms@k@species.arabiensis.wind_block_threshold_ms\endcsname{0.3375}
\expandafter\def\csname ms@k@species.arabiensis.wind_reference\endcsname{"u100/v100"}
\expandafter\def\csname ms@k@species.coluzzii.activity_weights\endcsname{"0.0, 0.255, 0.48, 0.265"}
\expandafter\def\csname ms@k@species.coluzzii.aestivation_supported\endcsname{true}
\expandafter\def\csname ms@k@species.coluzzii.anemotaxis_deviation_deg\endcsname{30.0}
\expandafter\def\csname ms@k@species.coluzzii.anemotaxis_rule\endcsname{"upwind"}
\expandafter\def\csname ms@k@species.coluzzii.arrested_daily_survival\endcsname{0.986}
\expandafter\def\csname ms@k@species.coluzzii.blood_meal_refractory_days\endcsname{0.5}
\expandafter\def\csname ms@k@species.coluzzii.blood_meal_size_ul\endcsname{0.3}
\expandafter\def\csname ms@k@species.coluzzii.blood_meal_target_ul\endcsname{0.5}
\expandafter\def\csname ms@k@species.coluzzii.casting_angle_deg\endcsname{90.0}
\expandafter\def\csname ms@k@species.coluzzii.casting_probability\endcsname{0.1}
\expandafter\def\csname ms@k@species.coluzzii.co2_activation_radius_m\endcsname{100.0}
\expandafter\def\csname ms@k@species.coluzzii.daily_move_capacity_m\endcsname{590.0}
\expandafter\def\csname ms@k@species.coluzzii.directed_flight_budget_h\endcsname{0.4856}
\expandafter\def\csname ms@k@species.coluzzii.eggs_per_ul\endcsname{223.0}
\expandafter\def\csname ms@k@species.coluzzii.exploration_range_m\endcsname{100.0}
\expandafter\def\csname ms@k@species.coluzzii.exploration_step_correlation\endcsname{0.5}
\expandafter\def\csname ms@k@species.coluzzii.exploration_step_m\endcsname{50.0}
\expandafter\def\csname ms@k@species.coluzzii.first_cycle_target_ul\endcsname{0.8}
\expandafter\def\csname ms@k@species.coluzzii.flight_speed_ms\endcsname{0.3375}
\expandafter\def\csname ms@k@species.coluzzii.gonotrophic_arrest_prob_dry\endcsname{0.2704}
\expandafter\def\csname ms@k@species.coluzzii.male_daily_survival\endcsname{0.773}
\expandafter\def\csname ms@k@species.coluzzii.male_ratio\endcsname{0.5}
\expandafter\def\csname ms@k@species.coluzzii.meal_size_min_ul\endcsname{0.1}
\expandafter\def\csname ms@k@species.coluzzii.olfactory_radius_m\endcsname{50.0}
\expandafter\def\csname ms@k@species.coluzzii.plume_memory_s\endcsname{0.7}
\expandafter\def\csname ms@k@species.coluzzii.pref.cattle\endcsname{0.005}
\expandafter\def\csname ms@k@species.coluzzii.pref.goat\endcsname{0.003}
\expandafter\def\csname ms@k@species.coluzzii.pref.human\endcsname{0.99}
\expandafter\def\csname ms@k@species.coluzzii.pref.sheep\endcsname{0.001}
\expandafter\def\csname ms@k@species.coluzzii.pref.wildlife\endcsname{0.001}
\expandafter\def\csname ms@k@species.coluzzii.probing_success_rate\endcsname{0.7275}
\expandafter\def\csname ms@k@species.coluzzii.salinity_hi_tol_ppt\endcsname{30.0}
\expandafter\def\csname ms@k@species.coluzzii.salinity_opt_ppt\endcsname{0.0}
\expandafter\def\csname ms@k@species.coluzzii.salinity_slope\endcsname{0.08}
\expandafter\def\csname ms@k@species.coluzzii.terminal_radius_m\endcsname{0.15}
\expandafter\def\csname ms@k@species.coluzzii.visual_radius_m\endcsname{0.485}
\expandafter\def\csname ms@k@species.coluzzii.weights.crop_pluvial\endcsname{0.05}
\expandafter\def\csname ms@k@species.coluzzii.weights.open_pluvial\endcsname{0.25}
\expandafter\def\csname ms@k@species.coluzzii.weights.surface_persistent\endcsname{0.4}
\expandafter\def\csname ms@k@species.coluzzii.weights.urban_persistent\endcsname{0.2}
\expandafter\def\csname ms@k@species.coluzzii.weights.wetland_persistent\endcsname{0.1}
\expandafter\def\csname ms@k@species.coluzzii.weights.wetland_pluvial\endcsname{0.0}
\expandafter\def\csname ms@k@species.coluzzii.wind_block_threshold_ms\endcsname{0.3375}
\expandafter\def\csname ms@k@species.coluzzii.wind_reference\endcsname{"u100/v100"}
\expandafter\def\csname ms@k@species.default_species\endcsname{"coluzzii"}
\expandafter\def\csname ms@k@species.funestus.activity_weights\endcsname{"0.02, 0.1, 0.8, 0.08"}
\expandafter\def\csname ms@k@species.funestus.aestivation_supported\endcsname{false}
\expandafter\def\csname ms@k@species.funestus.anemotaxis_deviation_deg\endcsname{30.0}
\expandafter\def\csname ms@k@species.funestus.anemotaxis_rule\endcsname{"upwind"}
\expandafter\def\csname ms@k@species.funestus.arrested_daily_survival\endcsname{0.986}
\expandafter\def\csname ms@k@species.funestus.blood_meal_refractory_days\endcsname{0.5}
\expandafter\def\csname ms@k@species.funestus.blood_meal_size_ul\endcsname{0.3}
\expandafter\def\csname ms@k@species.funestus.blood_meal_target_ul\endcsname{0.5}
\expandafter\def\csname ms@k@species.funestus.casting_angle_deg\endcsname{90.0}
\expandafter\def\csname ms@k@species.funestus.casting_probability\endcsname{0.1}
\expandafter\def\csname ms@k@species.funestus.co2_activation_radius_m\endcsname{100.0}
\expandafter\def\csname ms@k@species.funestus.daily_move_capacity_m\endcsname{590.0}
\expandafter\def\csname ms@k@species.funestus.directed_flight_budget_h\endcsname{0.4856}
\expandafter\def\csname ms@k@species.funestus.eggs_per_ul\endcsname{223.0}
\expandafter\def\csname ms@k@species.funestus.exploration_range_m\endcsname{100.0}
\expandafter\def\csname ms@k@species.funestus.exploration_step_correlation\endcsname{0.5}
\expandafter\def\csname ms@k@species.funestus.exploration_step_m\endcsname{50.0}
\expandafter\def\csname ms@k@species.funestus.first_cycle_target_ul\endcsname{0.8}
\expandafter\def\csname ms@k@species.funestus.flight_speed_ms\endcsname{0.3375}
\expandafter\def\csname ms@k@species.funestus.gonotrophic_arrest_prob_dry\endcsname{0.0}
\expandafter\def\csname ms@k@species.funestus.male_daily_survival\endcsname{0.773}
\expandafter\def\csname ms@k@species.funestus.male_ratio\endcsname{0.5}
\expandafter\def\csname ms@k@species.funestus.meal_size_min_ul\endcsname{0.1}
\expandafter\def\csname ms@k@species.funestus.olfactory_radius_m\endcsname{50.0}
\expandafter\def\csname ms@k@species.funestus.plume_memory_s\endcsname{0.7}
\expandafter\def\csname ms@k@species.funestus.pref.cattle\endcsname{0.005}
\expandafter\def\csname ms@k@species.funestus.pref.goat\endcsname{0.003}
\expandafter\def\csname ms@k@species.funestus.pref.human\endcsname{0.99}
\expandafter\def\csname ms@k@species.funestus.pref.sheep\endcsname{0.001}
\expandafter\def\csname ms@k@species.funestus.pref.wildlife\endcsname{0.001}
\expandafter\def\csname ms@k@species.funestus.probing_success_rate\endcsname{0.7275}
\expandafter\def\csname ms@k@species.funestus.salinity_hi_tol_ppt\endcsname{15.0}
\expandafter\def\csname ms@k@species.funestus.salinity_opt_ppt\endcsname{0.0}
\expandafter\def\csname ms@k@species.funestus.salinity_slope\endcsname{0.15}
\expandafter\def\csname ms@k@species.funestus.terminal_radius_m\endcsname{0.15}
\expandafter\def\csname ms@k@species.funestus.visual_radius_m\endcsname{0.485}
\expandafter\def\csname ms@k@species.funestus.weights.crop_pluvial\endcsname{0.05}
\expandafter\def\csname ms@k@species.funestus.weights.open_pluvial\endcsname{0.1}
\expandafter\def\csname ms@k@species.funestus.weights.surface_persistent\endcsname{0.45}
\expandafter\def\csname ms@k@species.funestus.weights.urban_persistent\endcsname{0.1}
\expandafter\def\csname ms@k@species.funestus.weights.wetland_persistent\endcsname{0.3}
\expandafter\def\csname ms@k@species.funestus.weights.wetland_pluvial\endcsname{0.0}
\expandafter\def\csname ms@k@species.funestus.wind_block_threshold_ms\endcsname{0.3375}
\expandafter\def\csname ms@k@species.funestus.wind_reference\endcsname{"u100/v100"}
\expandafter\def\csname ms@k@species.gambiae_ss.activity_weights\endcsname{"0.02, 0.1, 0.8, 0.08"}
\expandafter\def\csname ms@k@species.gambiae_ss.aestivation_supported\endcsname{true}
\expandafter\def\csname ms@k@species.gambiae_ss.anemotaxis_deviation_deg\endcsname{30.0}
\expandafter\def\csname ms@k@species.gambiae_ss.anemotaxis_rule\endcsname{"upwind"}
\expandafter\def\csname ms@k@species.gambiae_ss.arrested_daily_survival\endcsname{0.986}
\expandafter\def\csname ms@k@species.gambiae_ss.blood_meal_refractory_days\endcsname{0.5}
\expandafter\def\csname ms@k@species.gambiae_ss.blood_meal_size_ul\endcsname{0.3}
\expandafter\def\csname ms@k@species.gambiae_ss.blood_meal_target_ul\endcsname{0.5}
\expandafter\def\csname ms@k@species.gambiae_ss.casting_angle_deg\endcsname{90.0}
\expandafter\def\csname ms@k@species.gambiae_ss.casting_probability\endcsname{0.1}
\expandafter\def\csname ms@k@species.gambiae_ss.co2_activation_radius_m\endcsname{100.0}
\expandafter\def\csname ms@k@species.gambiae_ss.daily_move_capacity_m\endcsname{590.0}
\expandafter\def\csname ms@k@species.gambiae_ss.directed_flight_budget_h\endcsname{0.4856}
\expandafter\def\csname ms@k@species.gambiae_ss.eggs_per_ul\endcsname{223.0}
\expandafter\def\csname ms@k@species.gambiae_ss.exploration_range_m\endcsname{100.0}
\expandafter\def\csname ms@k@species.gambiae_ss.exploration_step_correlation\endcsname{0.5}
\expandafter\def\csname ms@k@species.gambiae_ss.exploration_step_m\endcsname{50.0}
\expandafter\def\csname ms@k@species.gambiae_ss.first_cycle_target_ul\endcsname{0.8}
\expandafter\def\csname ms@k@species.gambiae_ss.flight_speed_ms\endcsname{0.3375}
\expandafter\def\csname ms@k@species.gambiae_ss.gonotrophic_arrest_prob_dry\endcsname{0.6838}
\expandafter\def\csname ms@k@species.gambiae_ss.male_daily_survival\endcsname{0.773}
\expandafter\def\csname ms@k@species.gambiae_ss.male_ratio\endcsname{0.5}
\expandafter\def\csname ms@k@species.gambiae_ss.meal_size_min_ul\endcsname{0.1}
\expandafter\def\csname ms@k@species.gambiae_ss.olfactory_radius_m\endcsname{50.0}
\expandafter\def\csname ms@k@species.gambiae_ss.plume_memory_s\endcsname{0.7}
\expandafter\def\csname ms@k@species.gambiae_ss.pref.cattle\endcsname{0.005}
\expandafter\def\csname ms@k@species.gambiae_ss.pref.goat\endcsname{0.003}
\expandafter\def\csname ms@k@species.gambiae_ss.pref.human\endcsname{0.99}
\expandafter\def\csname ms@k@species.gambiae_ss.pref.sheep\endcsname{0.001}
\expandafter\def\csname ms@k@species.gambiae_ss.pref.wildlife\endcsname{0.001}
\expandafter\def\csname ms@k@species.gambiae_ss.probing_success_rate\endcsname{0.7275}
\expandafter\def\csname ms@k@species.gambiae_ss.salinity_hi_tol_ppt\endcsname{20.0}
\expandafter\def\csname ms@k@species.gambiae_ss.salinity_opt_ppt\endcsname{0.0}
\expandafter\def\csname ms@k@species.gambiae_ss.salinity_slope\endcsname{0.12}
\expandafter\def\csname ms@k@species.gambiae_ss.terminal_radius_m\endcsname{0.15}
\expandafter\def\csname ms@k@species.gambiae_ss.visual_radius_m\endcsname{0.485}
\expandafter\def\csname ms@k@species.gambiae_ss.weights.crop_pluvial\endcsname{0.1}
\expandafter\def\csname ms@k@species.gambiae_ss.weights.open_pluvial\endcsname{0.65}
\expandafter\def\csname ms@k@species.gambiae_ss.weights.surface_persistent\endcsname{0.05}
\expandafter\def\csname ms@k@species.gambiae_ss.weights.urban_persistent\endcsname{0.15}
\expandafter\def\csname ms@k@species.gambiae_ss.weights.wetland_persistent\endcsname{0.05}
\expandafter\def\csname ms@k@species.gambiae_ss.weights.wetland_pluvial\endcsname{0.0}
\expandafter\def\csname ms@k@species.gambiae_ss.wind_block_threshold_ms\endcsname{0.3375}
\expandafter\def\csname ms@k@species.gambiae_ss.wind_reference\endcsname{"u100/v100"}
\expandafter\def\csname ms@k@species.k_ref\endcsname{50.0}
\expandafter\def\csname ms@k@species.melas.activity_weights\endcsname{"0.02, 0.1, 0.8, 0.08"}
\expandafter\def\csname ms@k@species.melas.aestivation_supported\endcsname{false}
\expandafter\def\csname ms@k@species.melas.anemotaxis_deviation_deg\endcsname{30.0}
\expandafter\def\csname ms@k@species.melas.anemotaxis_rule\endcsname{"upwind"}
\expandafter\def\csname ms@k@species.melas.arrested_daily_survival\endcsname{0.986}
\expandafter\def\csname ms@k@species.melas.blood_meal_refractory_days\endcsname{0.5}
\expandafter\def\csname ms@k@species.melas.blood_meal_size_ul\endcsname{0.3}
\expandafter\def\csname ms@k@species.melas.blood_meal_target_ul\endcsname{0.5}
\expandafter\def\csname ms@k@species.melas.casting_angle_deg\endcsname{90.0}
\expandafter\def\csname ms@k@species.melas.casting_probability\endcsname{0.1}
\expandafter\def\csname ms@k@species.melas.co2_activation_radius_m\endcsname{100.0}
\expandafter\def\csname ms@k@species.melas.daily_move_capacity_m\endcsname{590.0}
\expandafter\def\csname ms@k@species.melas.directed_flight_budget_h\endcsname{0.4856}
\expandafter\def\csname ms@k@species.melas.eggs_per_ul\endcsname{223.0}
\expandafter\def\csname ms@k@species.melas.exploration_range_m\endcsname{100.0}
\expandafter\def\csname ms@k@species.melas.exploration_step_correlation\endcsname{0.5}
\expandafter\def\csname ms@k@species.melas.exploration_step_m\endcsname{50.0}
\expandafter\def\csname ms@k@species.melas.first_cycle_target_ul\endcsname{0.8}
\expandafter\def\csname ms@k@species.melas.flight_speed_ms\endcsname{0.3375}
\expandafter\def\csname ms@k@species.melas.gonotrophic_arrest_prob_dry\endcsname{0.0}
\expandafter\def\csname ms@k@species.melas.male_daily_survival\endcsname{0.773}
\expandafter\def\csname ms@k@species.melas.male_ratio\endcsname{0.5}
\expandafter\def\csname ms@k@species.melas.meal_size_min_ul\endcsname{0.1}
\expandafter\def\csname ms@k@species.melas.olfactory_radius_m\endcsname{50.0}
\expandafter\def\csname ms@k@species.melas.plume_memory_s\endcsname{0.7}
\expandafter\def\csname ms@k@species.melas.pref.cattle\endcsname{0.005}
\expandafter\def\csname ms@k@species.melas.pref.goat\endcsname{0.003}
\expandafter\def\csname ms@k@species.melas.pref.human\endcsname{0.99}
\expandafter\def\csname ms@k@species.melas.pref.sheep\endcsname{0.001}
\expandafter\def\csname ms@k@species.melas.pref.wildlife\endcsname{0.001}
\expandafter\def\csname ms@k@species.melas.probing_success_rate\endcsname{0.7275}
\expandafter\def\csname ms@k@species.melas.salinity_hi_tol_ppt\endcsname{38.0}
\expandafter\def\csname ms@k@species.melas.salinity_opt_ppt\endcsname{0.0}
\expandafter\def\csname ms@k@species.melas.salinity_slope\endcsname{0.02}
\expandafter\def\csname ms@k@species.melas.terminal_radius_m\endcsname{0.15}
\expandafter\def\csname ms@k@species.melas.visual_radius_m\endcsname{0.485}
\expandafter\def\csname ms@k@species.melas.weights.crop_pluvial\endcsname{0.05}
\expandafter\def\csname ms@k@species.melas.weights.open_pluvial\endcsname{0.4}
\expandafter\def\csname ms@k@species.melas.weights.surface_persistent\endcsname{0.25}
\expandafter\def\csname ms@k@species.melas.weights.urban_persistent\endcsname{0.15}
\expandafter\def\csname ms@k@species.melas.weights.wetland_persistent\endcsname{0.15}
\expandafter\def\csname ms@k@species.melas.weights.wetland_pluvial\endcsname{0.0}
\expandafter\def\csname ms@k@species.melas.wind_block_threshold_ms\endcsname{0.3375}
\expandafter\def\csname ms@k@species.melas.wind_reference\endcsname{"u100/v100"}
\expandafter\def\csname ms@k@species.neutral.weights.crop_pluvial\endcsname{0.05}
\expandafter\def\csname ms@k@species.neutral.weights.open_pluvial\endcsname{0.4}
\expandafter\def\csname ms@k@species.neutral.weights.surface_persistent\endcsname{0.25}
\expandafter\def\csname ms@k@species.neutral.weights.urban_persistent\endcsname{0.15}
\expandafter\def\csname ms@k@species.neutral.weights.wetland_persistent\endcsname{0.15}
\expandafter\def\csname ms@k@species.neutral.weights.wetland_pluvial\endcsname{0.0}
\expandafter\def\csname ms@k@species.stephensi.activity_weights\endcsname{"0.02, 0.1, 0.8, 0.08"}
\expandafter\def\csname ms@k@species.stephensi.aestivation_supported\endcsname{false}
\expandafter\def\csname ms@k@species.stephensi.anemotaxis_deviation_deg\endcsname{30.0}
\expandafter\def\csname ms@k@species.stephensi.anemotaxis_rule\endcsname{"upwind"}
\expandafter\def\csname ms@k@species.stephensi.arrested_daily_survival\endcsname{0.986}
\expandafter\def\csname ms@k@species.stephensi.blood_meal_refractory_days\endcsname{0.5}
\expandafter\def\csname ms@k@species.stephensi.blood_meal_size_ul\endcsname{0.3}
\expandafter\def\csname ms@k@species.stephensi.blood_meal_target_ul\endcsname{0.5}
\expandafter\def\csname ms@k@species.stephensi.casting_angle_deg\endcsname{90.0}
\expandafter\def\csname ms@k@species.stephensi.casting_probability\endcsname{0.1}
\expandafter\def\csname ms@k@species.stephensi.co2_activation_radius_m\endcsname{100.0}
\expandafter\def\csname ms@k@species.stephensi.daily_move_capacity_m\endcsname{590.0}
\expandafter\def\csname ms@k@species.stephensi.directed_flight_budget_h\endcsname{0.4856}
\expandafter\def\csname ms@k@species.stephensi.eggs_per_ul\endcsname{223.0}
\expandafter\def\csname ms@k@species.stephensi.exploration_range_m\endcsname{100.0}
\expandafter\def\csname ms@k@species.stephensi.exploration_step_correlation\endcsname{0.5}
\expandafter\def\csname ms@k@species.stephensi.exploration_step_m\endcsname{50.0}
\expandafter\def\csname ms@k@species.stephensi.first_cycle_target_ul\endcsname{0.8}
\expandafter\def\csname ms@k@species.stephensi.flight_speed_ms\endcsname{0.3375}
\expandafter\def\csname ms@k@species.stephensi.gonotrophic_arrest_prob_dry\endcsname{0.0}
\expandafter\def\csname ms@k@species.stephensi.male_daily_survival\endcsname{0.773}
\expandafter\def\csname ms@k@species.stephensi.male_ratio\endcsname{0.5}
\expandafter\def\csname ms@k@species.stephensi.meal_size_min_ul\endcsname{0.1}
\expandafter\def\csname ms@k@species.stephensi.olfactory_radius_m\endcsname{50.0}
\expandafter\def\csname ms@k@species.stephensi.plume_memory_s\endcsname{0.7}
\expandafter\def\csname ms@k@species.stephensi.pref.cattle\endcsname{0.15}
\expandafter\def\csname ms@k@species.stephensi.pref.goat\endcsname{0.05}
\expandafter\def\csname ms@k@species.stephensi.pref.human\endcsname{0.7}
\expandafter\def\csname ms@k@species.stephensi.pref.sheep\endcsname{0.05}
\expandafter\def\csname ms@k@species.stephensi.pref.wildlife\endcsname{0.05}
\expandafter\def\csname ms@k@species.stephensi.probing_success_rate\endcsname{0.7275}
\expandafter\def\csname ms@k@species.stephensi.salinity_hi_tol_ppt\endcsname{20.0}
\expandafter\def\csname ms@k@species.stephensi.salinity_opt_ppt\endcsname{0.0}
\expandafter\def\csname ms@k@species.stephensi.salinity_slope\endcsname{0.15}
\expandafter\def\csname ms@k@species.stephensi.terminal_radius_m\endcsname{0.15}
\expandafter\def\csname ms@k@species.stephensi.visual_radius_m\endcsname{0.485}
\expandafter\def\csname ms@k@species.stephensi.weights.crop_pluvial\endcsname{0.05}
\expandafter\def\csname ms@k@species.stephensi.weights.open_pluvial\endcsname{0.4}
\expandafter\def\csname ms@k@species.stephensi.weights.surface_persistent\endcsname{0.25}
\expandafter\def\csname ms@k@species.stephensi.weights.urban_persistent\endcsname{0.15}
\expandafter\def\csname ms@k@species.stephensi.weights.wetland_persistent\endcsname{0.15}
\expandafter\def\csname ms@k@species.stephensi.weights.wetland_pluvial\endcsname{0.0}
\expandafter\def\csname ms@k@species.stephensi.wind_block_threshold_ms\endcsname{0.3375}
\expandafter\def\csname ms@k@species.stephensi.wind_reference\endcsname{"u100/v100"}
\expandafter\def\csname ms@k@thermal.efd_max_eggs_per_day\endcsname{0.5}
\expandafter\def\csname ms@k@thermal.efd_tmax_c\endcsname{34.0}
\expandafter\def\csname ms@k@thermal.efd_tmin_c\endcsname{16.0}
\expandafter\def\csname ms@k@thermal.efd_topt_c\endcsname{25.0}
\expandafter\def\csname ms@k@thermal.egg_development_scale\endcsname{0.00062}
\expandafter\def\csname ms@k@thermal.egg_development_tmax_c\endcsname{40.0}
\expandafter\def\csname ms@k@thermal.egg_development_tmin_c\endcsname{10.0}
\expandafter\def\csname ms@k@thermal.larva_development_scale\endcsname{0.001}
\expandafter\def\csname ms@k@thermal.larva_development_tmax_c\endcsname{40.0}
\expandafter\def\csname ms@k@thermal.larva_development_tmin_c\endcsname{10.0}
\expandafter\def\csname ms@k@thermal.pupa_development_scale\endcsname{0.00066}
\expandafter\def\csname ms@k@thermal.pupa_development_tmax_c\endcsname{40.0}
\expandafter\def\csname ms@k@thermal.pupa_development_tmin_c\endcsname{12.0}
\expandafter\def\csname ms@k@transmission.beta_hv\endcsname{0.4}
\expandafter\def\csname ms@k@transmission.beta_vh\endcsname{0.5}
\expandafter\def\csname ms@k@transmission.human_immunity_days\endcsname{180}
\expandafter\def\csname ms@k@transmission.human_incubation_days\endcsname{12}
\expandafter\def\csname ms@k@transmission.human_infectious_days\endcsname{20}
\expandafter\def\csname ms@k@transmission.initial_human_prevalence\endcsname{0.05}
\expandafter\def\csname ms@k@transmission.initial_vector_infected_frac\endcsname{0.0}
\expandafter\def\csname ms@k@urban.beta_urban\endcsname{0.75}
\expandafter\def\csname ms@k@urban.rain_pools_built_damage\endcsname{0.3}
\expandafter\def\csname ms@k@version\endcsname{3}
\expandafter\def\csname ms@k@wildlife.class_suitability.bare\endcsname{0.2}
\expandafter\def\csname ms@k@wildlife.class_suitability.built\endcsname{0.05}
\expandafter\def\csname ms@k@wildlife.class_suitability.crops\endcsname{0.3}
\expandafter\def\csname ms@k@wildlife.class_suitability.flooded_vegetation\endcsname{0.45}
\expandafter\def\csname ms@k@wildlife.class_suitability.grass\endcsname{0.6}
\expandafter\def\csname ms@k@wildlife.class_suitability.shrub_and_scrub\endcsname{0.7}
\expandafter\def\csname ms@k@wildlife.class_suitability.snow_and_ice\endcsname{0.1}
\expandafter\def\csname ms@k@wildlife.class_suitability.trees\endcsname{0.9}
\expandafter\def\csname ms@k@wildlife.class_suitability.water\endcsname{0.0}
\expandafter\def\csname ms@k@wildlife.weights.habitat\endcsname{0.5}
\expandafter\def\csname ms@k@wildlife.weights.remoteness\endcsname{0.2}
\expandafter\def\csname ms@k@wildlife.weights.water\endcsname{0.3}

% Generated by tools/gen_concepts_values.py -- do not edit.
% Values are read from the CODE that implements each stage, by AST.
% A key that no longer exists makes \codeconst fail the LaTeX build
% with an error, in any document that references it.
% Each line carries the file:line the value was read from.

\makeatletter
\newcommand{\codeconst}[1]{%
  \expandafter\ifx\csname cc@k@#1\endcsname\relax
    \errmessage{Concept document: unknown code constant `#1'}%
  \else
    \csname cc@k@#1\endcsname%
  \fi}
\makeatother

% download.wildlife_habitat_fallback = 0.5  (mal-commonlib/src/mal_commonlib/data/loaders/wildlife.py:205)
\expandafter\def\csname cc@k@download.wildlife_habitat_fallback\endcsname{0.5}
% download.wildlife_remoteness_neutral = 1.0  (mal-commonlib/src/mal_commonlib/data/loaders/wildlife.py:161)
\expandafter\def\csname cc@k@download.wildlife_remoteness_neutral\endcsname{1.0}
% ingest.bbox_east = 1.5  (mal-commonlib/src/mal_commonlib/aois.yaml:18)
\expandafter\def\csname cc@k@ingest.bbox_east\endcsname{1.5}
% ingest.bbox_north = 11.5  (mal-commonlib/src/mal_commonlib/aois.yaml:18)
\expandafter\def\csname cc@k@ingest.bbox_north\endcsname{11.5}
% ingest.bbox_south = 4.5  (mal-commonlib/src/mal_commonlib/aois.yaml:18)
\expandafter\def\csname cc@k@ingest.bbox_south\endcsname{4.5}
% ingest.bbox_west = -3.5  (mal-commonlib/src/mal_commonlib/aois.yaml:18)
\expandafter\def\csname cc@k@ingest.bbox_west\endcsname{-3.5}
% ingest.cell_size_m = 1000  (mal-commonlib/src/mal_commonlib/aois.yaml:20)
\expandafter\def\csname cc@k@ingest.cell_size_m\endcsname{1000}
% ingest.env_daily_channels = 2  (mal-core/src/mal_core/ingest/daily_nc.py:414)
\expandafter\def\csname cc@k@ingest.env_daily_channels\endcsname{2}
% ingest.env_static_planes = 5  (mal-core/src/mal_core/ingest/daily_nc.py:414)
\expandafter\def\csname cc@k@ingest.env_static_planes\endcsname{5}
% ingest.fallback_building_rural = 0.1  (mal-core/src/mal_core/ingest/hosts.py:142)
\expandafter\def\csname cc@k@ingest.fallback_building_rural\endcsname{0.1}
% ingest.fallback_building_urban = 0.6  (mal-core/src/mal_core/ingest/hosts.py:142)
\expandafter\def\csname cc@k@ingest.fallback_building_urban\endcsname{0.6}
% ingest.fallback_wildlife_proxy = 0.3  (mal-core/src/mal_core/ingest/hosts.py:165)
\expandafter\def\csname cc@k@ingest.fallback_wildlife_proxy\endcsname{0.3}
% ingest.jrc_occurrence_threshold_pct = 80  (mal-commonlib/src/mal_commonlib/data/loaders/jrc_gsw.py:72)
\expandafter\def\csname cc@k@ingest.jrc_occurrence_threshold_pct\endcsname{80}
% ingest.k_ref_diag_per_km2 = 1329900.0  (mal-core/src/mal_core/habitat/capacity.py:30)
\expandafter\def\csname cc@k@ingest.k_ref_diag_per_km2\endcsname{1329900.0}
% ingest.mobility_beta_day = 0.05  (mal-core/src/mal_core/ingest/mobility.py:43)
\expandafter\def\csname cc@k@ingest.mobility_beta_day\endcsname{0.05}
% ingest.mobility_beta_livestock = 0.1  (mal-core/src/mal_core/ingest/mobility.py:45)
\expandafter\def\csname cc@k@ingest.mobility_beta_livestock\endcsname{0.1}
% ingest.mobility_beta_night = 0.5  (mal-core/src/mal_core/ingest/mobility.py:44)
\expandafter\def\csname cc@k@ingest.mobility_beta_night\endcsname{0.5}
% ingest.mobility_cell_size_km = 1.0  (mal-core/src/mal_core/ingest/mobility.py:42)
\expandafter\def\csname cc@k@ingest.mobility_cell_size_km\endcsname{1.0}
% ingest.mobility_max_distance_km = 50.0  (mal-core/src/mal_core/ingest/mobility.py:46)
\expandafter\def\csname cc@k@ingest.mobility_max_distance_km\endcsname{50.0}
% ingest.mobility_sparsity_threshold = 0.0005  (mal-core/src/mal_core/ingest/mobility.py:47)
\expandafter\def\csname cc@k@ingest.mobility_sparsity_threshold\endcsname{0.0005}
% ingest.urban_class_rural = 50  (mal-core/src/mal_core/ingest/hosts.py:126)
\expandafter\def\csname cc@k@ingest.urban_class_rural\endcsname{50}
% ingest.urban_class_urban = 30  (mal-core/src/mal_core/ingest/hosts.py:125)
\expandafter\def\csname cc@k@ingest.urban_class_urban\endcsname{30}
% ingest.urban_class_water = 20  (mal-core/src/mal_core/ingest/hosts.py:127)
\expandafter\def\csname cc@k@ingest.urban_class_water\endcsname{20}
% prediction.block_size = 128  (mal-core/src/mal_core/prediction/predictor.py:126)
\expandafter\def\csname cc@k@prediction.block_size\endcsname{128}
% prediction.contract_version_manifest = "1.1"  (mal-core/src/mal_core/prediction/registry.py:31)
\expandafter\def\csname cc@k@prediction.contract_version_manifest\endcsname{"1.1"}
% prediction.contract_version_sidecar = "1.1"  (mal-core/src/mal_core/prediction/predictor.py:146)
\expandafter\def\csname cc@k@prediction.contract_version_sidecar\endcsname{"1.1"}
% prediction.in_channels = 10  (mal-core/src/mal_core/prediction/registry.py:32)
\expandafter\def\csname cc@k@prediction.in_channels\endcsname{10}
% prediction.manifest_version = "1.0"  (mal-core/src/mal_core/prediction/registry.py:30)
\expandafter\def\csname cc@k@prediction.manifest_version\endcsname{"1.0"}
% prediction.out_channels = 6  (mal-core/src/mal_core/prediction/registry.py:33)
\expandafter\def\csname cc@k@prediction.out_channels\endcsname{6}
% prediction.scale_continental_m = 5000  (mal-core/src/mal_core/prediction/aggregator.py:29)
\expandafter\def\csname cc@k@prediction.scale_continental_m\endcsname{5000}
% prediction.scale_national_m = 2000  (mal-core/src/mal_core/prediction/aggregator.py:28)
\expandafter\def\csname cc@k@prediction.scale_national_m\endcsname{2000}
% prediction.scale_regional_m = 1000  (mal-core/src/mal_core/prediction/aggregator.py:27)
\expandafter\def\csname cc@k@prediction.scale_regional_m\endcsname{1000}
% scoring.d11_larva_band_max = 0.9  (mal-core/src/mal_core/scoring/scorers/d11_aquatic_structure.py:35)
\expandafter\def\csname cc@k@scoring.d11_larva_band_max\endcsname{0.9}
% scoring.d11_larva_band_min = 0.65  (mal-core/src/mal_core/scoring/scorers/d11_aquatic_structure.py:35)
\expandafter\def\csname cc@k@scoring.d11_larva_band_min\endcsname{0.65}
% scoring.d11_pupa_band_max = 0.15  (mal-core/src/mal_core/scoring/scorers/d11_aquatic_structure.py:36)
\expandafter\def\csname cc@k@scoring.d11_pupa_band_max\endcsname{0.15}
% scoring.d11_pupa_band_min = 0.02  (mal-core/src/mal_core/scoring/scorers/d11_aquatic_structure.py:36)
\expandafter\def\csname cc@k@scoring.d11_pupa_band_min\endcsname{0.02}
% scoring.d15_expansion_fraction = 0.2  (mal-core/src/mal_core/scoring/scorers/d15_persistence.py:125)
\expandafter\def\csname cc@k@scoring.d15_expansion_fraction\endcsname{0.2}
% scoring.d15_min_horizon_days = 30  (mal-core/src/mal_core/scoring/scorers/d15_persistence.py:97)
\expandafter\def\csname cc@k@scoring.d15_min_horizon_days\endcsname{30}
% scoring.d15_survival_fraction = 0.3  (mal-core/src/mal_core/scoring/scorers/d15_persistence.py:105)
\expandafter\def\csname cc@k@scoring.d15_survival_fraction\endcsname{0.3}
% scoring.d16_coverage_target = 0.5  (mal-core/src/mal_core/scoring/scorers/d16_detection_coverage.py:44)
\expandafter\def\csname cc@k@scoring.d16_coverage_target\endcsname{0.5}
% scoring.d16_radius_km_default = 5.0  (mal-core/src/mal_core/scoring/scorers/d16_detection_coverage.py:43)
\expandafter\def\csname cc@k@scoring.d16_radius_km_default\endcsname{5.0}
% scoring.d16_sensitivity_radius_far_km = 10.0  (mal-core/src/mal_core/scoring/scorers/d16_detection_coverage.py:45)
\expandafter\def\csname cc@k@scoring.d16_sensitivity_radius_far_km\endcsname{10.0}
% scoring.d16_sensitivity_radius_mid_km = 5.0  (mal-core/src/mal_core/scoring/scorers/d16_detection_coverage.py:45)
\expandafter\def\csname cc@k@scoring.d16_sensitivity_radius_mid_km\endcsname{5.0}
% scoring.d16_sensitivity_radius_near_km = 2.0  (mal-core/src/mal_core/scoring/scorers/d16_detection_coverage.py:45)
\expandafter\def\csname cc@k@scoring.d16_sensitivity_radius_near_km\endcsname{2.0}
% scoring.d19_activation_floor = 0.05  (mal-core/src/mal_core/scoring/scorers/d19_pool_persistence.py:81)
\expandafter\def\csname cc@k@scoring.d19_activation_floor\endcsname{0.05}
% scoring.d19_check_drying_weight = 0.2  (mal-core/src/mal_core/scoring/scorers/d19_pool_persistence.py:130)
\expandafter\def\csname cc@k@scoring.d19_check_drying_weight\endcsname{0.2}
% scoring.d19_check_floor_weight = 0.3  (mal-core/src/mal_core/scoring/scorers/d19_pool_persistence.py:134)
\expandafter\def\csname cc@k@scoring.d19_check_floor_weight\endcsname{0.3}
% scoring.d19_check_monotone_weight = 0.2  (mal-core/src/mal_core/scoring/scorers/d19_pool_persistence.py:132)
\expandafter\def\csname cc@k@scoring.d19_check_monotone_weight\endcsname{0.2}
% scoring.d19_check_persistence_weight = 0.3  (mal-core/src/mal_core/scoring/scorers/d19_pool_persistence.py:128)
\expandafter\def\csname cc@k@scoring.d19_check_persistence_weight\endcsname{0.3}
% scoring.d19_dry_days = 10  (mal-core/src/mal_core/scoring/scorers/d19_pool_persistence.py:84)
\expandafter\def\csname cc@k@scoring.d19_dry_days\endcsname{10}
% scoring.d19_pulse_days = 3  (mal-core/src/mal_core/scoring/scorers/d19_pool_persistence.py:78)
\expandafter\def\csname cc@k@scoring.d19_pulse_days\endcsname{3}
% scoring.d19_rain_pulse_mm = 30.0  (mal-core/src/mal_core/scoring/scorers/d19_pool_persistence.py:79)
\expandafter\def\csname cc@k@scoring.d19_rain_pulse_mm\endcsname{30.0}
% scoring.d1_band_max = 0.7  (mal-core/src/mal_core/scoring/scorers/d1_habitat_occupancy.py:33)
\expandafter\def\csname cc@k@scoring.d1_band_max\endcsname{0.7}
% scoring.d1_band_min = 0.2  (mal-core/src/mal_core/scoring/scorers/d1_habitat_occupancy.py:33)
\expandafter\def\csname cc@k@scoring.d1_band_min\endcsname{0.2}
% scoring.d1_center = 0.45  (mal-core/src/mal_core/scoring/scorers/d1_habitat_occupancy.py:34)
\expandafter\def\csname cc@k@scoring.d1_center\endcsname{0.45}
% scoring.d1_sigma = 0.15  (mal-core/src/mal_core/scoring/scorers/d1_habitat_occupancy.py:35)
\expandafter\def\csname cc@k@scoring.d1_sigma\endcsname{0.15}
% scoring.d20_cap_probe_mm = 80.0  (mal-core/src/mal_core/scoring/scorers/d20_washout_response.py:56)
\expandafter\def\csname cc@k@scoring.d20_cap_probe_mm\endcsname{80.0}
% scoring.d20_check_cap_weight = 0.2  (mal-core/src/mal_core/scoring/scorers/d20_washout_response.py:71)
\expandafter\def\csname cc@k@scoring.d20_check_cap_weight\endcsname{0.2}
% scoring.d20_check_no_washout_weight = 0.3  (mal-core/src/mal_core/scoring/scorers/d20_washout_response.py:69)
\expandafter\def\csname cc@k@scoring.d20_check_no_washout_weight\endcsname{0.3}
% scoring.d20_check_threshold_weight = 0.2  (mal-core/src/mal_core/scoring/scorers/d20_washout_response.py:73)
\expandafter\def\csname cc@k@scoring.d20_check_threshold_weight\endcsname{0.2}
% scoring.d20_check_washout_weight = 0.3  (mal-core/src/mal_core/scoring/scorers/d20_washout_response.py:67)
\expandafter\def\csname cc@k@scoring.d20_check_washout_weight\endcsname{0.3}
% scoring.d20_min_washout_fraction = 0.15  (mal-core/src/mal_core/scoring/scorers/d20_washout_response.py:48)
\expandafter\def\csname cc@k@scoring.d20_min_washout_fraction\endcsname{0.15}
% scoring.d20_no_washout_probe_mm = 30.0  (mal-core/src/mal_core/scoring/scorers/d20_washout_response.py:51)
\expandafter\def\csname cc@k@scoring.d20_no_washout_probe_mm\endcsname{30.0}
% scoring.d20_threshold_max_fraction = 0.01  (mal-core/src/mal_core/scoring/scorers/d20_washout_response.py:63)
\expandafter\def\csname cc@k@scoring.d20_threshold_max_fraction\endcsname{0.01}
% scoring.d20_threshold_probe_mm = 40.0  (mal-core/src/mal_core/scoring/scorers/d20_washout_response.py:61)
\expandafter\def\csname cc@k@scoring.d20_threshold_probe_mm\endcsname{40.0}
% scoring.d20_washout_probe_mm = 50.0  (mal-core/src/mal_core/scoring/scorers/d20_washout_response.py:46)
\expandafter\def\csname cc@k@scoring.d20_washout_probe_mm\endcsname{50.0}
% scoring.d21_horizon_day = 30  (mal-core/src/mal_core/scoring/scorers/d21_dispersal_kernel.py:112)
\expandafter\def\csname cc@k@scoring.d21_horizon_day\endcsname{30}
% scoring.d21_p50_band_max = 0.7  (mal-core/src/mal_core/scoring/scorers/d21_dispersal_kernel.py:39)
\expandafter\def\csname cc@k@scoring.d21_p50_band_max\endcsname{0.7}
% scoring.d21_p50_band_min = 0.15  (mal-core/src/mal_core/scoring/scorers/d21_dispersal_kernel.py:39)
\expandafter\def\csname cc@k@scoring.d21_p50_band_min\endcsname{0.15}
% scoring.d21_p50_center = 0.4  (mal-core/src/mal_core/scoring/scorers/d21_dispersal_kernel.py:41)
\expandafter\def\csname cc@k@scoring.d21_p50_center\endcsname{0.4}
% scoring.d21_p50_sigma = 0.2  (mal-core/src/mal_core/scoring/scorers/d21_dispersal_kernel.py:41)
\expandafter\def\csname cc@k@scoring.d21_p50_sigma\endcsname{0.2}
% scoring.d21_p90_band_max = 3.0  (mal-core/src/mal_core/scoring/scorers/d21_dispersal_kernel.py:40)
\expandafter\def\csname cc@k@scoring.d21_p90_band_max\endcsname{3.0}
% scoring.d21_p90_band_min = 0.7  (mal-core/src/mal_core/scoring/scorers/d21_dispersal_kernel.py:40)
\expandafter\def\csname cc@k@scoring.d21_p90_band_min\endcsname{0.7}
% scoring.d21_p90_center = 1.6  (mal-core/src/mal_core/scoring/scorers/d21_dispersal_kernel.py:42)
\expandafter\def\csname cc@k@scoring.d21_p90_center\endcsname{1.6}
% scoring.d21_p90_sigma = 0.8  (mal-core/src/mal_core/scoring/scorers/d21_dispersal_kernel.py:42)
\expandafter\def\csname cc@k@scoring.d21_p90_sigma\endcsname{0.8}
% scoring.d25_r_target = 0.6  (mal-core/src/mal_core/scoring/scorers/d25_cases_ghana.py:45)
\expandafter\def\csname cc@k@scoring.d25_r_target\endcsname{0.6}
% scoring.d2_lifetime_centre = 11.0  (mal-core/src/mal_core/scoring/scorers/d2_survival.py:70)
\expandafter\def\csname cc@k@scoring.d2_lifetime_centre\endcsname{11.0}
% scoring.d2_lifetime_width = 4.0  (mal-core/src/mal_core/scoring/scorers/d2_survival.py:70)
\expandafter\def\csname cc@k@scoring.d2_lifetime_width\endcsname{4.0}
% scoring.d2_window_max = 15  (mal-core/src/mal_core/scoring/scorers/d2_survival.py:47)
\expandafter\def\csname cc@k@scoring.d2_window_max\endcsname{15}
% scoring.d2_window_min = 8  (mal-core/src/mal_core/scoring/scorers/d2_survival.py:47)
\expandafter\def\csname cc@k@scoring.d2_window_min\endcsname{8}
% scoring.d3_band_max = 0.5  (mal-core/src/mal_core/scoring/scorers/d3_eip.py:69)
\expandafter\def\csname cc@k@scoring.d3_band_max\endcsname{0.5}
% scoring.d3_band_min = 0.2  (mal-core/src/mal_core/scoring/scorers/d3_eip.py:69)
\expandafter\def\csname cc@k@scoring.d3_band_min\endcsname{0.2}
% scoring.d3_center = 0.35  (mal-core/src/mal_core/scoring/scorers/d3_eip.py:72)
\expandafter\def\csname cc@k@scoring.d3_center\endcsname{0.35}
% scoring.d3_day = 30  (mal-core/src/mal_core/scoring/scorers/d3_eip.py:67)
\expandafter\def\csname cc@k@scoring.d3_day\endcsname{30}
% scoring.d3_sigma = 0.15  (mal-core/src/mal_core/scoring/scorers/d3_eip.py:72)
\expandafter\def\csname cc@k@scoring.d3_sigma\endcsname{0.15}
% scoring.d4_collapse_max = 0.8  (mal-core/src/mal_core/scoring/scorers/d4_stability.py:70)
\expandafter\def\csname cc@k@scoring.d4_collapse_max\endcsname{0.8}
% scoring.d4_collapse_min = 0.1  (mal-core/src/mal_core/scoring/scorers/d4_stability.py:70)
\expandafter\def\csname cc@k@scoring.d4_collapse_min\endcsname{0.1}
% scoring.d4_explosion_max = 5.0  (mal-core/src/mal_core/scoring/scorers/d4_stability.py:77)
\expandafter\def\csname cc@k@scoring.d4_explosion_max\endcsname{5.0}
% scoring.d5_band_max = 0.85  (mal-core/src/mal_core/scoring/scorers/d5_morans.py:95)
\expandafter\def\csname cc@k@scoring.d5_band_max\endcsname{0.85}
% scoring.d5_band_min = 0.2  (mal-core/src/mal_core/scoring/scorers/d5_morans.py:95)
\expandafter\def\csname cc@k@scoring.d5_band_min\endcsname{0.2}
% scoring.d5_center = 0.5  (mal-core/src/mal_core/scoring/scorers/d5_morans.py:98)
\expandafter\def\csname cc@k@scoring.d5_center\endcsname{0.5}
% scoring.d5_sigma = 0.25  (mal-core/src/mal_core/scoring/scorers/d5_morans.py:98)
\expandafter\def\csname cc@k@scoring.d5_sigma\endcsname{0.25}
% scoring.d8_min_score = 0.4  (mal-core/src/mal_core/scoring/scorers/d8_coupling.py:35)
\expandafter\def\csname cc@k@scoring.d8_min_score\endcsname{0.4}
% scoring.d8_pass_threshold = 0.3  (mal-core/src/mal_core/scoring/scorers/d8_coupling.py:76)
\expandafter\def\csname cc@k@scoring.d8_pass_threshold\endcsname{0.3}
% scoring.d8_score_divisor = 0.5  (mal-core/src/mal_core/scoring/scorers/d8_coupling.py:70)
\expandafter\def\csname cc@k@scoring.d8_score_divisor\endcsname{0.5}
% scoring.gate_census_reference = 27000000  (mal-core/src/mal_core/scoring/scorers/gates.py:191)
\expandafter\def\csname cc@k@scoring.gate_census_reference\endcsname{27000000}
% scoring.gate_ratio_band_max = 1.2  (mal-core/src/mal_core/scoring/scorers/gates.py:214)
\expandafter\def\csname cc@k@scoring.gate_ratio_band_max\endcsname{1.2}
% scoring.gate_ratio_band_min = 0.8  (mal-core/src/mal_core/scoring/scorers/gates.py:214)
\expandafter\def\csname cc@k@scoring.gate_ratio_band_min\endcsname{0.8}
% scoring.gate_row_sum_threshold = 0.99  (mal-core/src/mal_core/scoring/scorers/gates.py:271)
\expandafter\def\csname cc@k@scoring.gate_row_sum_threshold\endcsname{0.99}
% scoring.gate_row_sum_tolerance = 0.01  (mal-core/src/mal_core/scoring/scorers/gates.py:260)
\expandafter\def\csname cc@k@scoring.gate_row_sum_tolerance\endcsname{0.01}
% scoring.gate_rural_class = 50  (mal-core/src/mal_core/scoring/scorers/gates.py:317)
\expandafter\def\csname cc@k@scoring.gate_rural_class\endcsname{50}
% scoring.gate_sample_files = 48  (mal-core/src/mal_core/scoring/scorers/gates.py:59)
\expandafter\def\csname cc@k@scoring.gate_sample_files\endcsname{48}
% scoring.gate_urban_class = 30  (mal-core/src/mal_core/scoring/scorers/gates.py:316)
\expandafter\def\csname cc@k@scoring.gate_urban_class\endcsname{30}
% scoring.gate_urban_ratio_max = 0.7  (mal-core/src/mal_core/scoring/scorers/gates.py:336)
\expandafter\def\csname cc@k@scoring.gate_urban_ratio_max\endcsname{0.7}
% scoring.gate_wall_time_limit_s = 30.0  (mal-core/src/mal_core/scoring/scorers/gates.py:158)
\expandafter\def\csname cc@k@scoring.gate_wall_time_limit_s\endcsname{30.0}
% scoring.regression_threshold = -0.05  (mal-core/src/mal_core/scoring/feedback.py:51)
\expandafter\def\csname cc@k@scoring.regression_threshold\endcsname{-0.05}
% scoring.verdict_borderline = 0.5  (mal-core/src/mal_core/scoring/runner.py:24)
\expandafter\def\csname cc@k@scoring.verdict_borderline\endcsname{0.5}
% scoring.verdict_regressed = 0.3  (mal-core/src/mal_core/scoring/runner.py:24)
\expandafter\def\csname cc@k@scoring.verdict_regressed\endcsname{0.3}
% scoring.verdict_viable = 0.7  (mal-core/src/mal_core/scoring/runner.py:24)
\expandafter\def\csname cc@k@scoring.verdict_viable\endcsname{0.7}
% training.ckpt_period = 10  (mal-core/src/mal_core/training/trainer.py:102)
\expandafter\def\csname cc@k@training.ckpt_period\endcsname{10}
% training.dice_count_threshold = 0.5  (mal-core/src/mal_core/training/model.py:110)
\expandafter\def\csname cc@k@training.dice_count_threshold\endcsname{0.5}
% training.dice_eps = 1e-06  (mal-core/src/mal_core/training/model.py:97)
\expandafter\def\csname cc@k@training.dice_eps\endcsname{1e-06}
% training.dice_tau = 0.5  (mal-core/src/mal_core/training/model.py:87)
\expandafter\def\csname cc@k@training.dice_tau\endcsname{0.5}
% training.dice_thresh = 0.5  (mal-core/src/mal_core/training/model.py:87)
\expandafter\def\csname cc@k@training.dice_thresh\endcsname{0.5}
% training.dice_weight = 0.5  (mal-core/src/mal_core/training/model.py:105)
\expandafter\def\csname cc@k@training.dice_weight\endcsname{0.5}
% training.env_channels = 4  (mal-core/src/mal_core/training/dataset.py:22)
\expandafter\def\csname cc@k@training.env_channels\endcsname{4}
% training.fallback_in_channels = 6  (mal-core/src/mal_core/training/wrapper.py:62)
\expandafter\def\csname cc@k@training.fallback_in_channels\endcsname{6}
% training.fallback_out_channels = 2  (mal-core/src/mal_core/training/wrapper.py:66)
\expandafter\def\csname cc@k@training.fallback_out_channels\endcsname{2}
% training.pad_multiple = 16  (mal-core/src/mal_core/training/wrapper.py:100)
\expandafter\def\csname cc@k@training.pad_multiple\endcsname{16}
% training.patch_size = 128  (mal-core/src/mal_core/training/dataset.py:19)
\expandafter\def\csname cc@k@training.patch_size\endcsname{128}
% training.sub_seed = 42  (mal-core/src/mal_core/training/dataset.py:210)
\expandafter\def\csname cc@k@training.sub_seed\endcsname{42}
% training.transmission_channels = 4  (mal-core/src/mal_core/training/dataset.py:21)
\expandafter\def\csname cc@k@training.transmission_channels\endcsname{4}
% training.unet_channels = 32, 64, 128, 256  (mal-core/src/mal_core/training/model.py:46)
\expandafter\def\csname cc@k@training.unet_channels\endcsname{32, 64, 128, 256}
% training.vector_channels = 2  (mal-core/src/mal_core/training/dataset.py:20)
\expandafter\def\csname cc@k@training.vector_channels\endcsname{2}


\title{The \code{prediction} stage of the MalariaSentinel pipeline\\
\large A conceptual description}
\author{MalariaSentinel conceptual documentation layer}
\date{}

\begin{document}
\maketitle

\begin{abstract}
\noindent
This document explains the \code{prediction} stage: the sixth stage of the
MalariaSentinel pipeline, and the one that applies a trained surrogate of the
agent-based model (ABM) to a real area of interest (AOI) and writes a
multi-band GeoTIFF. It is written for a reader who has never opened the
code. It is not a specification and not a code reference: it names no
invariant, no call graph and no test. Every claim is cited to a file and a
line a reviewer can open, and every number that exists in the project's
configuration is emitted by the generator rather than typed by hand. The
document is deliberately blunt about the stage's limits, because the stage's
limits are large: it emulates one step of a model rather than forecasting
epidemiology, it degrades silently to zeros when its inputs are absent, and
several of its most consequential constants are literals in the source rather
than values in the frozen configuration.
\end{abstract}

\tableofcontents

\section{What this stage is for}

\subsection{Which source wins when two of them disagree}

Three authorities govern the statements in this document, and they are not
interchangeable. \emph{Values} --- the numbers the pipeline actually uses ---
come from \code{mal-core/config/model\_constants.yaml}, which reaches this
document through the generated \code{\textbackslash mskey} macros, and from the
code constants that the generator reads out of the implementation, which
reach it through the generated \code{\textbackslash codeconst} macros; where a
value is both configured and hard-coded in the code, the configured value
wins. \emph{Contract values and
contract names} --- version labels, band names, field names --- come from the
generated contract table (\code{docs/concepts/generated/contracts.tex}) and
from the generated code constants that carry the labels the implementation
writes, and that generated table wins over any prose that paraphrases it.
\emph{Semantics} --- what a step means, what its inputs are for, what its
outputs do and do not license a reader to conclude --- belong to this
document, and this document wins over a contract table or a configuration
file whenever the two disagree about meaning. When the three contradict one
another, the lower authority is the bug: the code is what runs, the
configuration is what is meant to govern the code, and this document is what
is meant to describe the code, so a contradiction means that one of the three
is wrong and has to be corrected rather than reconciled. The rule is stated
here because this document owns the semantics of the \code{prediction} stage;
there is no separate page that arbitrates between the three authorities, and
none is needed while each stage states the rule for the content it owns.

\subsection{The question it answers}

The pipeline as a whole exists to support malaria elimination programmes: it
ingests spatial and epidemiological data, simulates transmission, scores the
simulation against biological expectations, trains a surrogate on the
simulation's own output, and finally applies that surrogate to produce a
raster a decision-maker could look at
(\code{mal-core/src/mal\_core/prediction/predictor.py:1-9}, the module
docstring, which enumerates the stage's own five steps: load the model, load
the state and environment tensors, run inference, save the GeoTIFF, return
its path).

The question the \code{prediction} stage answers is therefore narrow and
mechanical: \emph{given an AOI, a scale, a date and a trained surrogate, what
does the surrogate say the next state of the system looks like, on that AOI's
grid?}

Concretely, the stage does exactly five things
(\code{mal-core/src/mal\_core/prediction/predictor.py:77-150}, the body of
\code{run\_prediction}, the single entry point): it resolves the AOI and its
scale-specific grid; it loads a checkpoint from the model registry; it loads
the ABM's state raster and the environment tensor for that AOI; it runs one
forward pass of the surrogate; and it aggregates the result to the requested
scale and writes it to disk as a GeoTIFF with a JSON sidecar. The path of the
written raster is returned to the caller
(\code{mal-core/src/mal\_core/prediction/predictor.py:150}).

What counts as having answered the question is therefore also mechanical: a
GeoTIFF exists at the returned path, it opens with a valid CRS and transform,
and its sidecar parses. That is the entirety of the stage's success
criterion as implemented; there is no skill score, no validation against
surveillance data, and no threshold that the output must clear
(\code{mal-core/src/mal\_core/prediction/predictor.py:89-132}, the inference
and write path; \code{mal-core/src/mal\_core/prediction/}, the package
contains no validation module of any kind).

\subsection{What it is explicitly not for}

The negative purpose is as important as the purpose, and the stage's own
code makes several of these exclusions visible.

\paragraph{It is not a measurement.}
The surrogate was fitted to the ABM's own snapshots, not to observations.
The training target is the \emph{next} ABM snapshot's bands
(\code{mal-core/src/mal\_core/training/dataset.py:242-252}), and the model's
header states the mapping it learns: state at time $t$ plus environment in,
state at time $t+1$ out
(\code{mal-core/src/mal\_core/training/model.py:1-6}). The stage's output is
a \emph{model of a model}. It is not an estimate of infection in human
populations, not a prevalence measurement, and not a probability of anything.

\paragraph{It is not an independent forecast.}
Because the surrogate reproduces one step of the simulator, the output
inherits every structural assumption of the simulator --- its habitat
representation, its mortality constants, its transmission model --- without
adding evidence. Where the ABM is wrong, the prediction is wrong in the same
direction, and it presents that error with the same confidence.

\paragraph{It does not train, and it does not simulate.}
Training lives in the \code{training} stage
(\code{mal-core/src/mal\_core/training/trainer.py:14-25}, the signature of
\code{train\_unet}); the ABM itself is the \code{abm} stage. The prediction
stage's own signature exposes only model selection, AOI, scale, date,
scenario, output directory and rollout directory
(\code{mal-core/src/mal\_core/prediction/predictor.py:66-76}); there is no
parameter that could cause it to train or to run the simulator.

\paragraph{It does not classify.}
The CLI help text for the \code{predict} subcommand says the stage will
``produce a classified risk map''
(\code{mal-core/src/mal\_core/cli.py:120-131}), but no classification is
performed anywhere. The network's output is written to the raster unchanged
(\code{mal-core/src/mal\_core/prediction/predictor.py:89} and
\code{mal-core/src/mal\_core/prediction/predictor.py:128}). There is no
threshold, no risk class, no probability scale, no uncertainty band and no
ensemble over seeds
(\code{mal-core/src/mal\_core/prediction/predictor.py:66-76}, the signature,
which has no parameter for any of these). The help text is aspirational; the
artefact is a surrogate state raster.

\paragraph{It does not model interventions.}
This is worth stating separately because the stage accepts a scenario and
converts it, then throws the conversion away. The interventions and climate
anomalies of a scenario are turned into a parameter dictionary
(\code{mal-core/src/mal\_core/prediction/predictor.py:86-87},
\code{params = interventions\_to\_params(scenario)}), and the dictionary is
never passed to the model: the inference call takes only the state and the
environment (\code{mal-core/src/mal\_core/prediction/predictor.py:89}). A
scenario therefore affects the output only through its \code{aoi},
\code{scale}, \code{year} and \code{month} fields, which the CLI copies over
the command-line arguments
(\code{mal-core/src/mal\_core/cli.py:133-140}). No intervention effect of any
kind is attributable to this stage's output.

\subsection{Where the stage sits in the pipeline}

The stage consumes two upstream artefacts and produces one:

\begin{center}
\begin{tabular}{@{}lll@{}}
\toprule
Producer & Artefact & Consumer here \\
\midrule
\code{abm} & state COG (per-day snapshots and/or the final state raster) &
state tensor \\
\code{abm} & transmission COG (optional) & state tensor, appended \\
\code{ingest} & environment NetCDF & environment tensor \\
\code{training} & checkpoint + \code{model.yaml} manifest & the surrogate \\
\midrule
\code{prediction} & risk GeoTIFF + JSON sidecar & the sidecar readers: the
\code{status} CLI command and the FastAPI endpoints \\
\bottomrule
\end{tabular}
\end{center}

The consumers inside the repository are the pair of sidecar readers: the
\code{status} CLI command
(\code{mal-core/src/mal\_core/cli.py:866-877}) and the two FastAPI endpoints
\code{GET /aoi/\{name\}/risk} and \code{GET /aoi/\{name\}/status}
(\code{mal-core/src/mal\_core/server.py:89-104}). No map, dashboard or
intervention-planning consumer exists anywhere in the tree.

\section{What the stage consumes}

This section is where the consumption is first described, so the band names,
the contract version and the resampling consequences all appear here rather
than only in the appendix.

\subsection{The ABM state raster}

\subsubsection{What it is and what it carries}

The stage consumes the \emph{state COG} written by the \code{abm} stage: a
two-band GeoTIFF on the ABM's own grid, with a JSON sidecar beside it. Under
the current state contract, version \mskey{contracts.cogs_state}, the two
bands are:

% Generated by tools/gen_concepts_values.py -- do not edit.
\paragraph{The bands the prediction consumes (the state contract)}
\begin{longtable}{@{}ll@{}}
\toprule
Band & Meaning \\
\midrule
\endhead
\code{adult\_patch\_count} & adults associated with their patch, post-dispersal \\
\code{adult\_cell\_count} & adults at their current cell \\
\bottomrule
\end{longtable}


Both bands are \emph{raw per-cell adult counts} --- numbers of individuals.
They are not densities, not normalised by any capacity constant, and not
bounded to the unit interval. The \code{abm} stage's own header states this
explicitly: the capacity normalisation was retired at contract 4.0, band 1
counts the adults associated with a cell through their patch after dispersal,
and band 2 counts the adults at their current position
(\code{mal-core/src/mal\_core/abm/include/mal\_abm\_fast/wire.hpp:124-133}).

The sidecar that accompanies the state COG carries a substantial amount of
metadata: the CRS, the affine transform, the AOI slug and scale, the year,
month and seed, the rollout count and index, the generator version, a hash of
the ABM parameters, the contract version, the band names, the NoData
sentinel and the raster shape
(\code{mal-core/src/mal\_core/abm/include/mal\_abm\_fast/output\_contract.hpp:36-56};
the full key list is reproduced in the generated contract table below).

% Generated by tools/gen_concepts_values.py -- do not edit.
\paragraph{The output contracts}
\begin{longtable}{@{}llllll@{}}
\toprule
Artefact & Producer & Consumers & Version & Bands / variables & Sentinel \\
\midrule
\endhead
\code{env} & \code{mal\_core.ingest.env.build\_env\_tensor} & abm, prediction, scoring & "3.0" & rainfall, air\_temp\_c, jrc\_persistent, lc\_built, lc\_wetland, lc\_open\_plus\_water, lc\_crop, land\_frac (optional), salinity\_ppt (optional) & none \\
\code{habitat} & \code{mal\_core.ingest.env.build\_env\_tensor} & abm &  & hab\_type, habitat\_potential, row, col & none \\
\code{state\_cog} & \code{mal\_abm\_fast::write\_state\_cog} & scoring, prediction, training & "4.0" & adult\_patch\_count, adult\_cell\_count & -9999.0 \\
\code{state\_sidecar} & \code{mal\_abm\_fast::write\_state\_sidecar} & scoring & "4.0" &  &  \\
\code{manifest} & \code{mal\_core.download.manifest.update\_dataset} & download, ingest, abm &  &  &  \\
\code{host\_static} & \code{mal\_core.ingest.hosts.build\_host\_dataset} & abm &  & humans, livestock, urban, buildings &  \\
\code{mobility\_day} & \code{mal\_core.ingest.mobility.build\_mobility\_dataset} & abm &  &  &  \\
\bottomrule
\end{longtable}

\paragraph{The state sidecar's keys, in emission order}
\begin{longtable}{@{}l@{}}
\toprule
Key \\
\midrule
\endhead
\code{crs} \\
\code{transform} \\
\code{aoi\_slug} \\
\code{scale} \\
\code{year} \\
\code{month} \\
\code{seed} \\
\code{n\_rollouts} \\
\code{rollout\_index} \\
\code{generator\_version} \\
\code{abm\_params\_hash} \\
\code{contract\_version} \\
\code{band\_names} \\
\code{nodata} \\
\code{shape} \\
\bottomrule
\end{longtable}


\subsubsection{The sidecar is not consulted}

The state loader reads the raster and nothing else. It opens the GeoTIFF
(\code{mal-core/src/mal\_core/prediction/state\_loader.py:82-86}), reads every
band, replaces cells equal to the file's declared NoData value with zero
(\code{mal-core/src/mal\_core/prediction/state\_loader.py:88-89}), and derives
the number of output channels from the number of bands in the file
(\code{mal-core/src/mal\_core/prediction/state\_loader.py:92}). The sidecar
JSON is never opened.

The consequences are worth stating as consequences, because they are the
kind that a reader could otherwise assume are handled:

\begin{itemize}
\item \textbf{Bands are read positionally.} The loader never looks up a band
by name. Band 1 of the file is taken to be band 1 of the tensor, band 2 to be
band 2. If a file's bands were in a different order, the loader would not
notice.
\item \textbf{No version gating.} Nothing in the stage compares the state
contract version in the sidecar against anything. A state raster produced
under a retired contract --- in particular under the pre-4.0 contract, whose
bands were normalised counts --- is consumed silently and produces
meaningless input to the surrogate. There is in fact no trained checkpoint
anywhere in the tree: the runs tree holds no \code{*.pt} and no
\code{model.yaml}, so nothing exists under the current contract and nothing
exists under a retired one either. The
redesign plan for the capacity retirement states that checkpoints trained
against the retired normalised bands are not reusable without retraining
(\code{docs/plans/completed/redesign/03-retirada-k-max-y-separacion-de-escalas.md:67-70}),
and the registry performs no such check
(\code{mal-core/src/mal\_core/prediction/registry.py:95-117}).
\item \textbf{No NoData harmonisation beyond the raster header.} The
sentinel used to blank cells is the one declared inside the GeoTIFF, not the
one in the sidecar; the two are the same value in practice, but the loader
only ever sees the former.
\end{itemize}

\subsubsection{How the snapshot for a date is chosen}

The stage is asked for a \code{month}. The loader converts that into a target
day: \code{target\_day = min(730, max(1, month * 30))}
(\code{mal-core/src/mal\_core/prediction/state\_loader.py:50}) --- the month
number times thirty, clamped into the range one to seven hundred and thirty
days. This is a coarse proxy for a calendar month, and it is the only place
where the stage's date argument reaches the data.

The candidate directories are searched in a fixed order
(\code{mal-core/src/mal\_core/prediction/state\_loader.py:38-47}):

\begin{enumerate}
\item an explicit directory, if the caller supplied one
  (\code{mal-core/src/mal\_core/prediction/predictor.py:75}, the
  \code{rollout\_dir} argument, exposed on the CLI as
  \code{--rollout-dir} at \code{mal-core/src/mal\_core/cli.py:116-119});
\item \code{runs/<aoi slug>};
\item \code{runs/abm};
\item \code{runs/1year-50r};
\item \code{runs/<aoi slug>\_2years};
\item \code{runs/test\_gif\_14d}.
\end{enumerate}

This list is hardcoded; it is not derived from a manifest or from any
inventory of the runs tree
(\code{mal-core/src/mal\_core/prediction/state\_loader.py:41-47}).

Within a directory, the loader first collects files whose names contain the
AOI slug and the string \code{day}, then all files matching \code{*\_day*.tif}
(\code{mal-core/src/mal\_core/prediction/state\_loader.py:57}); it removes
from that list any file whose name contains \code{transmission}
(\code{mal-core/src/mal\_core/prediction/state\_loader.py:59}). If nothing
remains, it falls back to files containing the slug and \code{seed}, then to
all TIFFs in the directory, again excluding transmission rasters
(\code{mal-core/src/mal\_core/prediction/state\_loader.py:61-62}). Among the
candidates it selects the one whose \code{\_dayNNN} suffix is numerically
closest to the target day; a file with no day suffix is treated as day
thirty (\code{mal-core/src/mal\_core/prediction/state\_loader.py:65-75}).

Two consequences follow directly from that loop's structure. First, the
search stops at the \emph{first} directory that yields any candidate
(\code{mal-core/src/mal\_core/prediction/state\_loader.py:76}, the
\code{break}), so a directory with a poor match beats a later directory with
a perfect one. Second, if no directory yields a candidate, the loader raises
\code{FileNotFoundError}
(\code{mal-core/src/mal\_core/prediction/state\_loader.py:78-79}) --- which the
caller converts into an all-zero state (see
Section~8).

\paragraph{Which files the ABM actually writes.}
The per-day snapshot file name is the output stem plus \code{\_dayNNN},
where \code{NNN} is zero-padded to three digits
(\code{mal-core/src/mal\_core/abm/src/main.cpp:169-183}), and the per-day
transmission file follows the same pattern
(\code{mal-core/src/mal\_core/abm/src/main.cpp:218-230}). Crucially, per-day
state snapshots are written only when the interval option is positive, and
that option's default is zero
(\code{mal-core/src/mal\_core/abm/src/main.cpp:1203-1216},
\code{mal-core/src/mal\_core/abm/src/main.cpp:409}), while the transmission
snapshot interval likewise defaults to zero, meaning the final transmission
snapshot only
(\code{mal-core/src/mal\_core/abm/src/main.cpp:806-809}). With default flags
the ABM writes only the final state raster and the final transmission raster.
The day-indexed files that the loader prefers are the exception, not the
norm, and a reader must not assume a month can be selected at all in a
default run.

\subsubsection{The transmission branch}

If the caller asks for transmission --- and the stage's only caller always
does (\code{mal-core/src/mal\_core/prediction/predictor.py:54},
\code{include\_transmission=True}) --- the loader looks for a transmission
raster in the directory of the chosen state file
(\code{mal-core/src/mal\_core/prediction/state\_loader.py:110-127}), trying
three strategies in order:

\begin{enumerate}
\item the paired name, i.e. the state file's stem plus \code{\_transmission.tif}
  (\code{mal-core/src/mal\_core/prediction/state\_loader.py:113-114});
\item any file matching \code{*transmission*dayNNN.tif}, where \code{NNN} is
  the day index parsed out of the chosen state file's name
  (\code{mal-core/src/mal\_core/prediction/state\_loader.py:116-122});
\item any file matching \code{*transmission*.tif}
  (\code{mal-core/src/mal\_core/prediction/state\_loader.py:124-127}).
\end{enumerate}

The transmission raster carries four bands, named
\code{human\_prevalence}, \code{human\_incidence},
\code{infectious\_vector\_pressure} and \code{active\_focus}
(\code{mal-core/src/mal\_core/abm/include/mal\_abm\_fast/transmission\_output.hpp:21-25}).
When one is found, its bands are resampled and concatenated onto the vector
bands, so the state tensor has as many channels as the state raster had bands
plus four --- six for the usual two-band vector raster, and more if the raster
carried more; when no transmission raster is found, the tensor has only the
vector bands
(\code{mal-core/src/mal\_core/prediction/state\_loader.py:36},
\code{mal-core/src/mal\_core/prediction/state\_loader.py:139-149}).

The third strategy is the one to be careful about. Because the default
transmission snapshot interval is zero
(\code{mal-core/src/mal\_core/abm/src/main.cpp:806-809}), the only transmission
raster in a default run is the \emph{final} one. The third strategy will
happily pick it up and pair it with a vector snapshot that is the nearest day
to the requested month --- two different points in time, concatenated without
comment. The stage does not record, anywhere, which day each channel came
from, and it does not refuse the pairing.

\subsubsection{Bring the state onto the AOI grid}

Both the state and the environment are brought onto the AOI's grid by
bilinear resampling (\code{mal-core/src/mal\_core/prediction/state\_loader.py:98-107};
the same operation appears in the environment loader at
\code{mal-core/src/mal\_core/prediction/env\_loader.py:74-83}). The target
grid is defined by the AOI's bounding box and by the cell count the AOI
object computes for the scale's resolution
(\code{mal-core/src/mal\_core/prediction/state\_loader.py:91-96}; the AOI's
own cell-count rule is a ceiling division of the bounding box's ground
dimensions by the resolution, at
\code{mal-commonlib/src/mal\_commonlib/aoi.py:195-199}). The source CRS is
taken from the source file and the destination CRS from the AOI
(\code{mal-core/src/mal\_core/prediction/state\_loader.py:99-107}).

The consequence of bilinear resampling a \emph{count-valued} band is worth
stating explicitly: after this step the values are no longer integers. A
cell that had two adults may become one and a half. The tensor the surrogate
receives is a field of fractional individuals, and the surrogate was trained
on integer-valued patches
(\code{mal-core/src/mal\_core/training/dataset.py:242-252}). Bilinear
interpolation also conserves nothing: it redistributes individuals between
neighbouring cells without preserving the total.

\subsection{The environment tensor}

\subsubsection{Which channels, and from where}

The environment tensor has four channels, in a fixed order
(\code{mal-core/src/mal\_core/prediction/env\_loader.py:17-18}):

\begin{center}
\begin{tabular}{@{}cl@{}}
\toprule
Index & Channel \\
\midrule
1 & \code{rainfall} \\
2 & \code{air\_temp\_c} \\
3 & \code{jrc\_persistent} \\
4 & \code{lc\_wetland} \\
\bottomrule
\end{tabular}
\end{center}

The producer is the \code{ingest} stage, which writes an environment NetCDF
under the environment contract version \mskey{contracts.env_nc}. That file
carries seven base variables: the two time series \code{rainfall} and
\code{air\_temp\_c}, and five static planes \code{jrc\_persistent},
\code{lc\_built}, \code{lc\_wetland}, \code{lc\_open\_plus\_water} and
\code{lc\_crop}
(\code{mal-core/src/mal\_core/ingest/daily\_nc.py:404-431}), plus up to two
optional ones that are added only when the corresponding input was available:
the static plane \code{land\_frac} and, when a salinity product was supplied,
the time series \code{salinity\_ppt}
(\code{mal-core/src/mal\_core/ingest/daily\_nc.py:443-457}); the file is
named \code{<aoi>\_regional\_<y0>\_<y1>\_env.nc}
(\code{mal-core/src/mal\_core/ingest/daily\_nc.py:466}). The prediction stage
reads four of them: the two time series, \code{jrc\_persistent} and
\code{lc\_wetland}
(\code{mal-core/src/mal\_core/prediction/env\_loader.py:17-18}).

\subsubsection{The fallback chain}

The environment loader resolves its input through a chain of five
possibilities, and every step of the chain is a choice a reader can vary
(\code{mal-core/src/mal\_core/prediction/env\_loader.py:36-151}):

\begin{enumerate}
\item an explicit path, if one is supplied
  (\code{mal-core/src/mal\_core/prediction/env\_loader.py:36-39});
\item otherwise, a NetCDF in the AOI's data directory, matched by
  \code{*\_env.nc} and then by \code{*.nc}, taking the first in sorted order
  (\code{mal-core/src/mal\_core/prediction/env\_loader.py:41-47});
\item otherwise, a GeoTIFF in the same directory, matched by \code{*\_env*.tif}
  and then \code{*.tif} (\code{mal-core/src/mal\_core/prediction/env\_loader.py:45},
  \code{mal-core/src/mal\_core/prediction/env\_loader.py:48-49});
\item otherwise, the legacy array file \code{runs/env\_stack.npz}
  (\code{mal-core/src/mal\_core/prediction/env\_loader.py:120-148});
\item otherwise, an all-zero tensor
  (\code{mal-core/src/mal\_core/prediction/env\_loader.py:150-151}).
\end{enumerate}

Two properties of this chain deserve comment.

First, the data directory is addressed relatively, as
\code{Path("data") / aoi.slug}
(\code{mal-core/src/mal\_core/prediction/env\_loader.py:42}), so the stage
finds the AOI's data only if it is run from the repository root. The
convention elsewhere in the project is that the manifest at
\code{data/<aoi>/manifest.json} is the single source of truth for an AOI's
data inventory, and the prediction package contains a resolver that implements
that convention
(\code{mal-core/src/mal\_core/prediction/aoi\_resolver.py:68-100}), but that
resolver has no caller anywhere in the package: the environment loader
bypasses it with its own glob. This is a departure from the project's
manifest-first rule, and it is a departure that is present in the code, not a
missing feature that a reader might imagine.

Second, the NetCDF branch selects only the variables that are actually
present in the file, in the order of the channel list, and zero-fills any
that are missing
(\code{mal-core/src/mal\_core/prediction/env\_loader.py:56-66}). A file that
lacks \code{lc\_wetland} therefore yields a fourth channel of zeros --- a
silently different tensor with the same shape, not an error.

\subsubsection{The environment is not date-specific}

For a variable with a time axis, the loader takes the \emph{first} time step
of the file and discards the rest: \code{vals = vals[0]}
(\code{mal-core/src/mal\_core/prediction/env\_loader.py:60-62}). The stage's
\code{month} argument selects a state snapshot but does not select an
environment slice. The environmental forcing the surrogate sees is therefore
the first day of the ingested record, whatever month was requested.

This is a behavioural fact, not an intention stated in the code, and it is
listed among the open questions of the stage's question set. A reader who
wants to know what the surrogate actually saw must know it: the environment
channels do not correspond to the same date as the state channels.

\subsubsection{Bring the environment onto the AOI grid}

The NetCDF branch assumes that the file covers exactly the AOI's bounding
box: it reconstructs the source transform from the AOI's own bounds and the
file's own array dimensions
(\code{mal-core/src/mal\_core/prediction/env\_loader.py:70-73}), then
bilinearly resamples into the AOI grid
(\code{mal-core/src/mal\_core/prediction/env\_loader.py:74-83}). The GeoTIFF
branch instead takes the source CRS and transform from the file
(\code{mal-core/src/mal\_core/prediction/env\_loader.py:92-95}). The legacy
array branch reconstructs its transform from a bounds quadruple stored in the
file, falling back to a hardcoded default when the file does not carry one
(\code{mal-core/src/mal\_core/prediction/env\_loader.py:134-135}); that default
is a literal in the source and is listed in the evidence appendix.

\section{Choosing and loading a checkpoint}

\subsection{The registry and its manifests}

The model registry (\code{mal-core/src/mal\_core/prediction/registry.py}) is
a directory scanner. Each model is described by a \code{model.yaml} manifest
sitting beside its checkpoint, and the registry's manifest record
(\code{ModelManifest}, \code{mal-core/src/mal\_core/prediction/registry.py:24-31})
has the following fields:

\begin{center}
\begin{tabular}{@{}lll@{}}
\toprule
Field & Default & Meaning \\
\midrule
\code{name} & --- & the model's key \\
\code{version} & \code{\codeconst{prediction.manifest_version}} & the version component of the key \\
\code{contract\_version} & \code{\codeconst{prediction.contract_version_manifest}} & the model's own contract label \\
\code{in\_channels} & \code{\codeconst{prediction.in_channels}} & channels the network expects on input \\
\code{out\_channels} & \code{\codeconst{prediction.out_channels}} & channels the network emits \\
\code{checkpoint} & empty & file name of the weights, relative to the manifest \\
\code{description} & empty & free text \\
\bottomrule
\end{tabular}
\end{center}

The defaults are placeholders rather than measurements: a manifest that omits
the channel counts advertises a \codeconst{prediction.in_channels}-channel
input and a \codeconst{prediction.out_channels}-channel output
regardless of what the weights actually are. That disagreement is not
silent: the wrapper builds the network with the counts it was given and then
loads the checkpoint into it
(\code{mal-core/src/mal\_core/training/wrapper.py:64-65}), and
\code{load\_state\_dict} raises on any shape mismatch, so a manifest whose
counts disagree with its checkpoint cannot be loaded at all --- the run
aborts with the exception instead of running a mis-sized network. The
manifest's
\code{contract\_version} is a label on the model; it is not the state
contract, it is never compared with the state contract, and nothing in the
repository validates it.

The scan walks two roots --- the registry's own models directory, by default
\code{runs/models}, and the runs directory itself --- and collects every
\code{model.yaml} found recursively
(\code{mal-core/src/mal\_core/prediction/registry.py:53-72}). Each manifest
is keyed by its name and version joined with an at-sign. A manifest that
fails to parse is skipped silently
(\code{mal-core/src/mal\_core/prediction/registry.py:69-70}).

\subsection{Version resolution}

When a version is given, the key is the exact name-at-version pair
(\code{mal-core/src/mal\_core/prediction/registry.py:78-79}). When no version
is given, the registry takes the \emph{last} key in scan order that begins
with the requested name and an at-sign
(\code{mal-core/src/mal\_core/prediction/registry.py:81-89}); failing that, it
tries the bare name as a key
(\code{mal-core/src/mal\_core/prediction/registry.py:83-85}), and failing
that it raises a \code{KeyError}
(\code{mal-core/src/mal\_core/prediction/registry.py:87}).

``Last in scan order'' is not a semantic-version comparison, but it is
deterministic. The scan visits the manifests in sorted path order
(\code{mal-core/src/mal\_core/prediction/registry.py:61}) and the keys are
inserted in that order, so the ``last'' match is the manifest whose path is
lexicographically greatest among those carrying the requested name. Nothing
compares version strings: which of two versions of the same model wins is
decided by their directory paths, not by which is newer. The
CLI exposes no way to ask for a specific version: the \code{predict}
subcommand has a \code{--model} option and no version option
(\code{mal-core/src/mal\_core/cli.py:107-119}).

\subsection{Loading the weights}

The loader has a special case first: the literal name \code{"dummy"}
short-circuits to a zero model, bypassing the registry entirely
(\code{mal-core/src/mal\_core/prediction/registry.py:95-97}). Otherwise it
resolves the manifest, and if the manifest's checkpoint field is empty it
also returns the zero model
(\code{mal-core/src/mal\_core/prediction/registry.py:115-116}). A non-empty
checkpoint is resolved relative to the manifest's directory; if that file
does not exist, the loader raises \code{FileNotFoundError}
(\code{mal-core/src/mal\_core/prediction/registry.py:103-109}). The exception
is not caught: the registry is loaded outside any \code{try} block
(\code{mal-core/src/mal\_core/prediction/predictor.py:81}), and the two
exception types the stage treats as ``no input'' are the state loader's
\code{FileNotFoundError} and the environment loader's catch-all
(\code{mal-core/src/mal\_core/prediction/predictor.py:55-62}; see
Section~8). A missing checkpoint therefore aborts the run rather than
degrading to zeros.

The weights are loaded by a wrapper (\code{UNetWrapper},
\code{mal-core/src/mal\_core/training/wrapper.py:18-106}) that implements the
same prediction interface the registry expects
(\code{ModelProtocol}, \code{mal-core/src/mal\_core/prediction/registry.py:20-21}).
The wrapper resolves the channel counts in this order:

\begin{enumerate}
\item the counts the caller passed, which for the registry path are the
  manifest's counts
  (\code{mal-core/src/mal\_core/prediction/registry.py:110-114});
\item otherwise, a \code{model.yaml} sitting beside the checkpoint file,
  which supplies both counts, but only when the caller supplied none: the
  manifest is read as a fallback and the caller's value wins if it is set
  (\code{mal-core/src/mal\_core/training/wrapper.py:39-48}). On the registry
  path the counts are always supplied
  (\code{mal-core/src/mal\_core/prediction/registry.py:110-114}), so the
  manifest beside the checkpoint is never read;
\item otherwise, counts inferred from the weight tensors themselves: the
  input count from the first encoder convolution's second dimension, and the
  output count from the head convolution's first dimension
  (\code{mal-core/src/mal\_core/training/wrapper.py:50-59}).
\end{enumerate}

The weights are then loaded into the network and the network is put into
evaluation mode (\code{mal-core/src/mal\_core/training/wrapper.py:64-66}).
The compute device is chosen in the order CUDA, then Apple's Metal
accelerator, then CPU
(\code{mal-core/src/mal\_core/training/wrapper.py:29-37}).

\subsection{The dummy model is not a model}

The zero model returned for the name \code{"dummy"}, and for any manifest
with an empty checkpoint, returns a tensor of zeros with the same number of
channels as the state it was given and the same spatial dimensions
(\code{mal-core/src/mal\_core/prediction/registry.py:41-45}). Its output band
count therefore follows the \emph{input's} channel count, not the manifest's
declared output count: a two-channel state yields a two-band raster, a
six-channel state a six-band raster. It performs no arithmetic, and it
therefore exercises the file-writing path of the stage without exercising the
network at all.

\subsection{Who writes the manifests}

The \code{training} stage does. After training, the trainer writes a
\code{model.yaml} beside the weights, with the name taken from the output
directory's name, version \code{\codeconst{prediction.manifest_version}},
contract \code{\codeconst{prediction.contract_version_manifest}}, the
checkpoint file name \code{best\_model.pt}, and the input and output channel
counts \emph{measured} from the first training sample rather than declared
(\code{mal-core/src/mal\_core/training/trainer.py:95-106};
the measurement is at \code{mal-core/src/mal\_core/training/trainer.py:50-52}).
That measurement is the reason a trained model's manifest can be trusted even
though the schema's defaults are placeholders.

\section{Inference}

\subsection{What is executed}

One forward pass of the U-Net over the whole AOI grid, and nothing else. The
network (\code{UNet}, \code{mal-core/src/mal\_core/training/model.py:28-60})
is a four-stage encoder-decoder with skip connections, whose head is a
one-by-one convolution applied directly to the decoded features
(\code{mal-core/src/mal\_core/training/model.py:46}); there is no activation
on the output. The sigmoid appears only inside the loss function, as a
soft-thresholding device for the Dice term
(\code{mal-core/src/mal\_core/training/model.py:63-73}), and therefore plays
no part in the values written to disk.

\subsection{Assembling the input tensor}

The wrapper concatenates the state and the environment along the channel axis
(\code{mal-core/src/mal\_core/training/wrapper.py:79-88}). If the
concatenation has fewer channels than the network expects, zero channels are
appended; if it has more, it is truncated. The raster is then padded up to a
multiple of \codeconst{training.pad_multiple} rows and columns --- the network pools four times --- by
edge replication, and the result is cropped back to the original dimensions
after the forward pass
(\code{mal-core/src/mal\_core/training/wrapper.py:90-104}).

The channel order is the concatenation of the state channels and the
environment channels. The state's own order agrees with the training order:
the vector bands first, then the transmission bands
(\code{mal-core/src/mal\_core/prediction/state\_loader.py:149}, compared with
\code{mal-core/src/mal\_core/training/dataset.py:242-252}). The environment's
order does \emph{not} agree with the training order: the training dataset
selects the first four variables of its own list, which are
\code{rainfall}, \code{air\_temp\_c}, \code{jrc\_persistent} and
\code{lc\_built}
(\code{mal-core/src/mal\_core/training/dataset.py:72}), whereas the
prediction loader selects \code{rainfall}, \code{air\_temp\_c},
\code{jrc\_persistent} and \code{lc\_wetland}
(\code{mal-core/src/mal\_core/prediction/env\_loader.py:18}). A checkpoint
trained under the training order receives, at inference, a fourth channel
that means something different from the one it was trained on. The stage
contains no list that both sides read, and the mismatch is silent.

\subsection{Consequences of the padding rules}

Two consequences of the wrapper's rules are worth drawing out, because they
change the meaning of the output rather than merely its shape.

First, with the manifest's default input count of ten and the two most
common input shapes, the wrapper's padding rules bite exactly once. A
six-channel state (two vector bands plus four transmission bands) plus a
four-channel environment is exactly ten channels and needs no padding. A
two-channel state --- the case in which no transmission raster was found ---
plus four environment channels is six, so four zero channels are appended ---
and they are appended at the \emph{end} of the concatenation
(\code{mal-core/src/mal\_core/training/wrapper.py:80-88}). The layout the
network therefore receives is \code{[vector(2), env(4), zeros(4)]}, whereas
the training layout is
\code{[vector(2), transmission(4), env(4)]}
(\code{mal-core/src/mal\_core/training/dataset.py:242-243},
\code{mal-core/src/mal\_core/training/dataset.py:251}). The zero padding
does not land in the transmission positions: the \emph{environment} does. The
consequence is a misalignment rather than a shortage of data: with a state
file that carries no transmission raster, the network receives the environment
channels where it was trained to expect transmission channels, and zeros where
it was trained to expect the environment. The output band count is the
manifest's declared output count, so the stage still writes
\codeconst{prediction.out_channels} bands and
names them after the training target's order
(\code{mal-core/src/mal\_core/prediction/predictor.py:33-40}); nothing in the
code ties those names to what the network actually produced.

Second, a truncation would occur if the inputs exceeded the network's expected
count, in which case the trailing environment channels would be cut. With the
current defaults and shapes this cannot happen, but the rule is unconditional
in the wrapper.

\subsection{What the numbers mean}

The written values are raw network outputs. Their \emph{intended} units are the
units of the training target: for the first two bands, adult counts in the
sense of contract \mskey{contracts.cogs_state}; for the remaining bands, the
transmission quantities the surrogate was trained to reproduce. That mapping is
an intention, not a guarantee the code enforces: the head is a bare
convolution with no output activation
(\code{mal-core/src/mal\_core/training/model.py:46}), so the values are
whatever the mean-squared-error fit produced, and nothing constrains them to
be counts rather than, say, thousands of counts. They are not probabilities,
not percentages, and not clamped to any interval: the stage writes the
activations unchanged
(\code{mal-core/src/mal\_core/prediction/predictor.py:128}).

\section{Aggregation by scale}

\subsection{The scales and their grids}

The stage supports three scales, each with its own grid resolution and its
own aggregator (\code{mal-core/src/mal\_core/prediction/aggregator.py:25-29},
\code{mal-core/src/mal\_core/prediction/aggregator.py:74-79}). The AOI object
is built from the registry entry for the slug, with the registry's own
resolution overridden by the scale's
(\code{mal-core/src/mal\_core/prediction/aggregator.py:32-44}). The
resolution values are literals in the source; the generator now reads them
out of the code and emits them as the three per-scale resolution macros used
in the table below, and the evidence appendix cites their
source locations.

\begin{center}
\begin{tabular}{@{}llll@{}}
\toprule
Scale & Resolution & Grid cell count & Aggregation \\
\midrule
regional & \codeconst{prediction.scale_regional_m}~m & $(H, W)$ & identity \\
national & \codeconst{prediction.scale_national_m}~m & $(H, W)$ & rectangular block mean \\
continental & \codeconst{prediction.scale_continental_m}~m & $1 \times 1$ & single scalar \\
\bottomrule
\end{tabular}
\end{center}

\subsection{What each aggregator does}

\paragraph{Regional.} The raster is returned unchanged
(\code{mal-core/src/mal\_core/prediction/aggregator.py:51-53}). This is the
only scale for which the written raster's shape equals the grid the
transform was computed on.

\paragraph{National.} The output is a block mean: the raster is trimmed to a
whole number of blocks, reshaped into blocks, and averaged within each block
(\code{mal-core/src/mal\_core/prediction/aggregator.py:56-66}). The number
of block rows is the raster height divided by ten, at least one, and likewise
for the columns, so the effective block size is approximately ten by ten
cells of the \codeconst{prediction.scale_national_m}-metre grid. Each output
cell is therefore the mean over roughly a twenty-kilometre square.

This is a rectangular block mean on the grid, and it is worth naming the
difference plainly, because the grid definition's own description string says
something else: it describes the national scale as a mean-pool to the
second administrative level
(\code{mal-core/src/mal\_core/prediction/aggregator.py:27}). The code
pools into rectangles; it does not intersect administrative boundaries, and
no boundary geometry is loaded anywhere in the package. The description
string is not what the code does.

\paragraph{Continental.} The whole raster is reduced to one number, its
mean over all cells and all bands
(\code{mal-core/src/mal\_core/prediction/aggregator.py:69-71}). The writer
then reshapes a one-dimensional result into a single-pixel square raster
(\code{mal-core/src/mal\_core/prediction/predictor.py:102-108}), so the
continental output is a one-by-one, one-band raster --- a single cell, not a
grid --- whose value is the mean
of everything. Because its band count is one, the band-name table gives it
the name \code{predicted\_risk}
(\code{mal-core/src/mal\_core/prediction/predictor.py:30-31}) --- a name that
is applied on the basis of the band count alone, and that describes the
value no more accurately than the other scales' names describe theirs.

\subsection{Georeferencing does not follow the aggregation}

The transform written into the output is computed from the \emph{scale grid}
dimensions, before the aggregated array's shape is known
(\code{mal-core/src/mal\_core/prediction/predictor.py:78},
\code{mal-core/src/mal\_core/prediction/predictor.py:100}). The raster's
height and width are taken from the \emph{aggregated} array
(\code{mal-core/src/mal\_core/prediction/predictor.py:112},
\code{mal-core/src/mal\_core/prediction/predictor.py:117-118}). For the
regional scale the two coincide. For the national scale the written raster
has one tenth of the rows and columns of the grid the transform was built
for, while the transform's pixel size is still the
\codeconst{prediction.scale_national_m}-metre grid's pixel
size; the result is a raster whose pixels are twice the intended width and
whose extent covers roughly one tenth of the AOI's bounding box, anchored at
its upper-left corner. For the continental scale the same discrepancy is
extreme: a single pixel carrying a transform whose pixel size is the
\codeconst{prediction.scale_continental_m}-metre grid's, so the one-pixel
raster covers a small corner of the
bounding box rather than the whole of it.

This is stated here as a description of current behaviour, not as an
endorsement. Whether it is a defect or the intended contract is an open
question in the stage's question set; a reader who needs the national or
continental output in a map should not assume it is georeferenced to the AOI.

\section{What the stage produces}

\subsection{File name and location}

By default the output directory is the \code{predictions} directory inside the
runs directory (\code{mal-core/src/mal\_core/prediction/predictor.py:94-95});
a caller may override it. The file name is built from the AOI slug, the
scale, the year and the zero-padded month, followed by the suffix
\code{\_risk.tif}
(\code{mal-core/src/mal\_core/prediction/predictor.py:97}). The scale is
embedded in the name, which is what makes the ``latest prediction'' lookup
behave in the way described in Section~6.5.

\subsection{Raster format}

\begin{center}
\begin{tabular}{@{}ll@{}}
\toprule
Property & Value \\
\midrule
Driver & GeoTIFF \\
Data type & 32-bit floating point \\
Band count & the aggregated array's channel count \\
Tiling & tiled, blocks of \codeconst{prediction.block_size} by \codeconst{prediction.block_size} \\
Compression & deflate \\
CRS & the AOI's CRS \\
Transform & affine transform from the AOI bounds over the scale grid \\
NoData & the ingest sentinel, \mskey{ingest.nodata_sentinel} \\
\bottomrule
\end{tabular}
\end{center}

All of these are set in one profile dictionary
(\code{mal-core/src/mal\_core/prediction/predictor.py:113-126}). The NoData
value is read from the frozen configuration rather than typed in
(\code{mal-core/src/mal\_core/prediction/predictor.py:121}, reading the
\code{ingest.nodata\_sentinel} key), which is the stage's only numeric
constant with a configured home.

The NoData value is, in practice, never encountered in the data: the
surrogate's outputs are unbounded network activations, and the writer performs
no masking
(\code{mal-core/src/mal\_core/prediction/predictor.py:128}). The sentinel in
the profile is a declaration of what the file's \emph{absence} would look
like, not a description of any value present in it.

\subsection{Band names}

The band descriptions are chosen by band count, from a table
(\code{mal-core/src/mal\_core/prediction/predictor.py:30-41}):

\begin{center}
\begin{tabular}{@{}llp{7cm}@{}}
\toprule
Band count & Names & Applies to \\
\midrule
1 & \code{predicted\_risk} & the continental scalar \\
2 & \code{predicted\_adult\_patch\_count}, \code{predicted\_adult\_cell\_count}
  & a state without transmission \\
6 & the two adult-count names, then \code{predicted\_human\_prevalence},
  \code{predicted\_human\_incidence}, \code{predicted\_infectious\_pressure},
  \code{predicted\_transmission\_focus} & the full state \\
other & \code{band\_1} \dots & fallback \\
\bottomrule
\end{tabular}
\end{center}

Any other band count gets generic names
(\code{mal-core/src/mal\_core/prediction/predictor.py:129}). Note that the
six-band names mirror the surrogate's \emph{training target} order --- two
adult-count bands then four transmission bands --- and not the ABM's own band
names, which are \code{adult\_patch\_count} and \code{adult\_cell\_count}
(\code{mal-core/src/mal\_core/abm/include/mal\_abm\_fast/wire.hpp:124-133}).
The \code{predicted\_} prefix is the only thing marking the difference
between a simulated quantity and a surrogate's estimate of it, and it lives
in the band description rather than in the data.

\subsection{The sidecar}

Beside the raster, the stage writes a JSON file whose name is the raster's
name with the suffix \code{.tif.json}
(\code{mal-core/src/mal\_core/prediction/predictor.py:147}). The suffix is
the opposite of the convention used by the existing artefacts in the runs
tree, where a sidecar \emph{replaces} the raster's extension: the ABM's writer
turns \code{<name>.tif} into \code{<name>.json}
(\code{mal-core/src/mal\_core/abm/src/output\_contract.cpp:43-64},
\code{mal-core/src/mal\_core/abm/src/output\_contract.cpp:90-94}). The
prediction stage appends instead, so the two sidecars of the same stem can
coexist. The sidecar contains:

\begin{center}
\begin{tabular}{@{}ll@{}}
\toprule
Key & Meaning \\
\midrule
\code{crs} & the AOI's CRS \\
\code{transform} & the first six components of the affine transform \\
\code{aoi\_slug} & the requested slug \\
\code{scale} & the requested scale \\
\code{year} & the requested year \\
\code{month} & the requested month \\
\code{bands} & the band names written \\
\code{model\_name} & the requested model name \\
\code{model\_version} & the requested version, or the word \code{latest} \\
\code{contract\_version} & the output contract label \\
\code{generator\_version} & the writing tool's version \\
\bottomrule
\end{tabular}
\end{center}

All of these are written in one dictionary
(\code{mal-core/src/mal\_core/prediction/predictor.py:134-148}).

Two facts about the sidecar must be stated together, because the second
undercuts the first. First, the sidecar does not record \emph{which model was
actually loaded}. The registry may have resolved the requested name to some
other name-at-version pair --- the version resolution rule is ``last in scan
order'' --- but the sidecar writes the \code{model\_name} it was given and
the \code{model\_version} it was given, which the CLI never supplies, so the
version field always reads \code{latest}
(\code{mal-core/src/mal\_core/prediction/predictor.py:142-143};
the CLI's absence of a version option is at
\code{mal-core/src/mal\_core/cli.py:107-119}). There is no checkpoint hash,
no reference to the training run, and no fingerprint of the inputs.

Second, the sidecar's \code{contract\_version} is the stage's \emph{own}
output contract label. It is not the state contract version
\mskey{contracts.cogs_state} of the raster it consumed, and the two are
independent: the stage does not gate on the state contract at all, and no
consumer in the repository validates the prediction's own contract label.
The label and the generator version are literals in the source
(\code{mal-core/src/mal\_core/prediction/predictor.py:144-145}) rather than
values in the frozen configuration. The label is emitted by the generator as
\codeconst{prediction.contract_version_sidecar}; the generator version has no
generated home and is listed in the evidence appendix.

\subsection{The ``latest prediction''}
\label{sec:latest}

The lookup for an AOI's most recent prediction sorts the matching file names
and takes the last
(\code{mal-core/src/mal\_core/prediction/predictor.py:153-158}). The match
pattern is \code{<slug>\_*\_risk.tif}. Because the scale is embedded in the
name, the sort is lexicographic over a mixed set of scales and dates, so it
is neither chronological nor a single-scale series: \code{national} sorts
after \code{continental} and before \code{regional}, and the year-month
component only breaks ties within a scale. The reader of this fact should
treat the returned path as ``some prediction for this AOI'', not as ``the
most recent one''.

\section{The order of operations, and what follows from it}

\subsection{Pseudocode}

In the listing below, \code{PAD\_MULTIPLE} stands for the wrapper's padding
multiple, \codeconst{training.pad_multiple}.

\begin{verbatim}
resolve AOI and scale grid
load checkpoint (registry scan -> name@version -> weights -> device)
load state tensor   (search runs tree -> pick snapshot -> resample
                      -> optionally concatenate transmission bands)
load env tensor     (NetCDF -> GeoTIFF -> legacy array -> zeros -> resample)
assemble input      (concat(state, env), zero-pad or truncate to in_channels,
                      pad to a multiple of PAD_MULTIPLE)
forward pass        (one pass of the U-Net; crop back to original size)
aggregate           (identity | block mean | scalar)
compute transform   (from the scale grid, before the aggregate is inspected)
reshape             (scalar -> 1x1; 1-D -> square; 2-D -> single band)
write GeoTIFF       (bands, descriptions, profile)
write sidecar       (metadata dictionary)
return the raster path
\end{verbatim}

\subsection{Schedule table}

\begin{center}
\begin{longtable}{@{}llll@{}}
\toprule
Step & Actor & Order & Mode \\
\midrule
\endhead
Resolve AOI, scale grid & aggregator & first & deterministic \\
Load checkpoint & registry, wrapper & second & may raise \\
Load state tensor & state loader & third & may degrade to zeros \\
Load env tensor & environment loader & fourth & may degrade to zeros \\
Assemble input & wrapper & fifth & deterministic \\
Forward pass & network & sixth & deterministic \\
Aggregate & aggregator & seventh & deterministic \\
Compute transform & predictor & eighth & deterministic \\
Write raster & predictor & ninth & deterministic \\
Write sidecar & predictor & tenth & deterministic \\
\bottomrule
\end{longtable}
\end{center}

\subsection{Consequences of this order}

\begin{enumerate}
\item \textbf{The transform precedes the aggregation.} The georeferencing is
fixed from the scale grid before the aggregated array exists
(\code{mal-core/src/mal\_core/prediction/predictor.py:100}), while the
raster's dimensions come from the aggregated array afterwards
(\code{mal-core/src/mal\_core/prediction/predictor.py:112}). For the
regional scale the two agree; for the other two they do not, as described in
Section~5.3.
\item \textbf{Aggregation precedes writing, and aggregation changes the shape
that is written.} The band count in the file is the aggregated array's
channel count, not the network's output count in general: the continental
aggregator collapses all channels into one, so a six-band network output
becomes a one-band file, and the band-name table consequently assigns that
band the name \code{predicted\_risk}
(\code{mal-core/src/mal\_core/prediction/predictor.py:129}).
\item \textbf{The model is loaded before the inputs, and the inputs cannot
influence the choice of model.} Because the registry resolves the model
before any data is read, a run with missing inputs and a run with present
inputs use the same weights.
\item \textbf{The inputs are loaded before the output directory is created.}
The output directory is created only after inference
(\code{mal-core/src/mal\_core/prediction/predictor.py:94-95}), so a failed
inference leaves no partial output. Conversely, a successful inference
followed by a write failure leaves a raster without a sidecar.
\item \textbf{The scenario's parameters are computed and discarded before the
forward pass} (\code{mal-core/src/mal\_core/prediction/predictor.py:86-89}),
which is why no intervention can influence the output.
\item \textbf{Only the month reaches the data.} The year selects no
snapshot, no environment slice and no model; it appears in the file name and
the sidecar and nowhere else
(\code{mal-core/src/mal\_core/prediction/predictor.py:69-70},
\code{mal-core/src/mal\_core/prediction/predictor.py:97},
\code{mal-core/src/mal\_core/prediction/predictor.py:140-141}).
\end{enumerate}

\section{Declaration of absences}

\subsection{Absent inputs degrade silently}
\label{sec:absent-inputs}

This is the single most consequential fact in the stage. The input loader
catches two failure classes and converts both into all-zero tensors, without
a warning, without a log line and without any record in the output
(\code{mal-core/src/mal\_core/prediction/predictor.py:50-62}):

\begin{itemize}
\item a \code{FileNotFoundError} from the state loader --- meaning no ABM
  snapshot was found in any candidate directory --- yields a zero state of
  six channels;
\item \emph{any} exception from the environment loader, of whatever kind,
  yields a zero environment.
\end{itemize}

The stage then proceeds exactly as it would have: it runs the surrogate on
zeros, aggregates, and writes a complete, well-formed GeoTIFF with a
well-formed sidecar. Nothing in the output distinguishes this artefact from
one computed on real data. A consumer reading the sidecar sees a valid
contract version and a valid generator version; a consumer reading the raster
sees numbers.

Two narrower points about the same mechanism. The state branch catches only
\code{FileNotFoundError}; a corrupt or unreadable state file raises whatever
exception the reader produced, and that exception propagates and aborts the
run (\code{mal-core/src/mal\_core/prediction/predictor.py:55-56}). The
environment branch catches everything, so a corrupt environment file is
indistinguishable from an absent one. And because the state fallback is
hardcoded to six channels
(\code{mal-core/src/mal\_core/prediction/predictor.py:56}), the zero state
always has the full channel count, so the downstream zero-padding rule never
fires for a missing state --- the surrogate receives zeros in the
transmission channels, as it would in the two-channel case.

\subsection{Implemented but never reached}

\begin{itemize}
\item \textbf{The dummy input generator.} The helper that fabricates a
  zero state and a zero environment
  (\code{mal-core/src/mal\_core/prediction/predictor.py:44-47}) has no
  caller anywhere in the repository.
\item \textbf{The manifest-aware AOI resolver.} The resolver that reads
  \code{data/<aoi>/manifest.json} and returns the AOI's data files
  (\code{mal-core/src/mal\_core/prediction/aoi\_resolver.py:68-100}) is
  never called. The environment loader uses its own glob instead
  (\code{mal-core/src/mal\_core/prediction/env\_loader.py:42-49}), which is
  a departure from the project's manifest-first convention.
\end{itemize}

\subsection{Inputs accepted and not used}

\begin{itemize}
\item \textbf{The seed.} The state loader's signature accepts a seed,
  described in its docstring as the rollout seed to use
  (\code{mal-core/src/mal\_core/prediction/state\_loader.py:23},
  \code{mal-core/src/mal\_core/prediction/state\_loader.py:32}), and the
  body never refers to it. The seed is inert.
\item \textbf{The year.} The year selects no snapshot, no environment and no
  model. It is written into the file name and the sidecar.
\end{itemize}

\subsection{States the code never enters}

\begin{itemize}
\item The transmission branch is skipped entirely when no transmission raster
  is found (\code{mal-core/src/mal\_core/prediction/state\_loader.py:129-151});
  in that case the state tensor has two channels and the wrapper appends
  four zero channels.
\item The scenario branch computes a parameter dictionary and discards it
  (\code{mal-core/src/mal\_core/prediction/predictor.py:86-87}); the forward
  pass is called without it.
\item The \code{"dummy"} name and the empty-checkpoint case both return the
  zero model without touching the registry's manifests
  (\code{mal-core/src/mal\_core/prediction/registry.py:96-97},
  \code{mal-core/src/mal\_core/prediction/registry.py:115-116}).
\end{itemize}

\subsection{Unknown AOIs fail loudly}

The one failure mode in the stage that is \emph{not} silent is an unknown AOI
slug. The AOI is built from the registry of AOIs, and a slug that is not in
that registry raises a \code{ValueError} naming the slug and listing the
available ones
(\code{mal-commonlib/src/mal\_commonlib/aoi.py:275-285}; the call is at
\code{mal-core/src/mal\_core/prediction/aggregator.py:34}). The stage
therefore fails fast on an unknown area, in contrast with its silent
degradation on missing data.

\subsection{Whole capabilities that do not exist}

There is no ensemble over seeds, no uncertainty quantification, no
calibration of the output against any reference, no thresholding or
classification, no re-projection onto administrative boundaries, and no
epidemiological validation of any kind
(\code{mal-core/src/mal\_core/prediction/predictor.py:66-76}, the signature;
and the package contains no validation module). The stage writes what the
network produced.

\section{The seams with the other stages}

\subsection{With the \code{abm} stage}

The seam is the state COG and its sidecar. The prediction stage reads the
raster positionally and ignores the sidecar, so the seam is one-way and
unguarded: the ABM writes a rich description of what it produced, and the
prediction stage takes only the numbers. The ABM's own contract version
\mskey{contracts.cogs_state} is therefore not enforced at this seam; the
consequence of a mismatch is a surrogate fed normalised counts where it
expects raw ones, and no error.

The second half of the seam is the transmission raster, which may come from a
different day than the vector snapshot, as described in
Section~2.1.4.

\subsection{With the \code{training} stage}

The seam is the checkpoint and its manifest. The training stage measures the
network's input and output channel counts from its own data and writes them
into the manifest
(\code{mal-core/src/mal\_core/training/trainer.py:50-52},
\code{mal-core/src/mal\_core/training/trainer.py:95-106}); the prediction
stage reads them and uses them to size the network
(\code{mal-core/src/mal\_core/prediction/registry.py:110-114}).

Two mismatches live on this seam, and neither is detected.

\begin{itemize}
\item \textbf{The environment channel order.} Training and inference read
  different fourth channels, as described in Section~4.2.
\item \textbf{The environment may be empty at training and real at
  inference.} The training dataset substitutes an all-zero environment patch
  whenever no environment file was supplied
  (\code{mal-core/src/mal\_core/training/dataset.py:246-249}), and the
  training CLI's environment-path option defaults to nothing
  (\code{mal-core/src/mal\_core/cli.py:780}). A checkpoint trained that way
  has never seen a real environment, and the prediction stage will hand it
  one.
\end{itemize}

Neither the manifest's contract label nor any other field closes these gaps.

\section{What the artefact must not be used for}

\begin{enumerate}
\item \textbf{Not as a measurement of infection.} The output is a surrogate's
  estimate of what the ABM would do next; the ABM itself is a model, and
  the surrogate is a model of that model
  (\code{mal-core/src/mal\_core/training/model.py:1-6}).
\item \textbf{Not as a probability.} The values are unbounded network
  outputs, not a score in the unit interval
  (\code{mal-core/src/mal\_core/training/model.py:46},
  \code{mal-core/src/mal\_core/prediction/predictor.py:128}).
\item \textbf{Not as a risk class.} Despite the file name's suffix and the
  band name \code{predicted\_risk}, no classification is performed.
\item \textbf{Not as the result of an intervention scenario.} Interventions
  are converted and discarded
  (\code{mal-core/src/mal\_core/prediction/predictor.py:86-87}).
\item \textbf{Not as a dated product.} The state channels come from the
  snapshot nearest the requested month; the environment channels come from
  the first time step of the ingested file
  (\code{mal-core/src/mal\_core/prediction/env\_loader.py:60-62}); and the
  transmission channels may come from a different day again
  (\code{mal-core/src/mal\_core/prediction/state\_loader.py:124-127}).
\item \textbf{Not as a record of what was run.} The sidecar names the model
  that was requested, not the model that was loaded, and does not record
  that the inputs were absent
  (\code{mal-core/src/mal\_core/prediction/predictor.py:142-143},
  \code{mal-core/src/mal\_core/prediction/predictor.py:55-61}).
\item \textbf{Not as a national or continental map.} Their georeferencing
  does not match their dimensions
  (\code{mal-core/src/mal\_core/prediction/predictor.py:100},
  \code{mal-core/src/mal\_core/prediction/predictor.py:117-118}).
\end{enumerate}

\section{Evidence appendix}

\subsection{Implementation notes for a replicator}

\begin{itemize}
\item The state loader's directory search stops at the first directory that
  has any candidate, not the best match
  (\code{mal-core/src/mal\_core/prediction/state\_loader.py:52-76}). A
  replicator who puts a better snapshot in a later directory will not reach
  it.
\item A snapshot file with no \code{\_dayNNN} suffix is treated as day
  thirty when the nearest-day comparison is made
  (\code{mal-core/src/mal\_core/prediction/state\_loader.py:70}).
\item The environment loader's data directory is relative to the process's
  working directory (\code{mal-core/src/mal\_core/prediction/env\_loader.py:42}),
  so the stage must be run from the repository root to find the AOI's data.
\item The environment loader's NetCDF branch assumes the file spans exactly
  the AOI's bounding box (\code{mal-core/src/mal\_core/prediction/env\_loader.py:70-73}).
\item The wrapper's zero-padding and truncation rules are unconditional
  (\code{mal-core/src/mal\_core/training/wrapper.py:80-88}); the replicator
  should check the channel arithmetic before assuming a missing band is
  absent rather than zeroed.
\item The registry's version resolution is lexicographic over manifest paths,
  not over version strings
  (\code{mal-core/src/mal\_core/prediction/registry.py:61},
  \code{mal-core/src/mal\_core/prediction/registry.py:81-89}); a replicator
  who needs a specific version must edit the call, since the command line
  offers no version option (\code{mal-core/src/mal\_core/cli.py:107-119}).
\item The whole array is written to the file in a single call, and only the
  band descriptions are set afterwards, one band at a time
  (\code{mal-core/src/mal\_core/prediction/predictor.py:127-132}); a reader
  that treats a missing description as an error would misread the
  fallback names.
\item The sidecar's transform is truncated to its first six components
  (\code{mal-core/src/mal\_core/prediction/predictor.py:136}), which is the
  affine tuple without the offset convention spelled out.
\end{itemize}

\subsection{Values that are literals in the source, not in the configuration}

The project's documentation layer requires that values come from the frozen
configuration or from a generated table. Most of the values this stage uses
now have a generated home: the contract labels, the channel counts, the
per-scale resolutions, the padding multiple and the tiling block size are
read out of the implementation by the generator and reach this document as
\code{\textbackslash codeconst} macros. The values listed below are the ones that are still
literals in the source with no generated home --- neither a key in
\code{mal-core/config/model\_constants.yaml} nor a constant the generator
reads. They are listed here as a reading of the code, with their locations,
so that a reader can check them and so that the remaining gap is visible.

\begin{center}
\begin{longtable}{@{}llll@{}}
\toprule
Value & Where & Kind & Note \\
\midrule
\endhead
\code{"mal-execution-0.1.0"} (generator version) &
  \code{mal-core/src/mal\_core/prediction/predictor.py:145} &
  literal string & written into the sidecar \\
10 (block divisor, national) &
  \code{mal-core/src/mal\_core/prediction/aggregator.py:59} &
  literal integer & rows and columns per block \\
$(-3.0, 4.0, 2.0, 12.0)$ (legacy array bounds) &
  \code{mal-core/src/mal\_core/prediction/env\_loader.py:134} &
  literal tuple & fallback bounds, west/south/east/north \\
30 (days per month proxy) &
  \code{mal-core/src/mal\_core/prediction/state\_loader.py:50} &
  literal integer & date proxy \\
1, 730 (day clamp) &
  \code{mal-core/src/mal\_core/prediction/state\_loader.py:50} &
  literal integers & date clamp \\
\bottomrule
\end{longtable}
\end{center}

The legacy bounds are the fallback the environment loader uses when a legacy
array file carries no \code{bounds} entry of its own
(\code{mal-core/src/mal\_core/prediction/env\_loader.py:134}); they are not
the ingest bounding box of any AOI, and they are not read from
\code{mal-commonlib/src/mal\_commonlib/aois.yaml}. They describe a box
shifted north-east relative to the Ghana AOI's own bounds (west $-3.5$,
south $4.5$, east $1.5$, north $11.5$), reaching further east and north than
the AOI does, so a legacy array file without bounds is georeferenced onto a
different area than the one the rest of the stage works on.

The configuration file currently defines two contract versions, the
environment contract and the state contract
(\code{mal-core/config/model\_constants.yaml:8-10}); there is no key for the
prediction stage's own output contract, and no row for the prediction
artefact in the generated contract table
(\code{docs/concepts/generated/contracts.tex}). The stage's contract label
is nevertheless emitted, as a code constant read out of the implementation
(\codeconst{prediction.contract_version_sidecar}), as are the per-scale
resolutions; what remains without a generated home is the set of values
listed in the table above. These are gaps rather than facts, and the
question set that accompanies this document records them as open.

\subsection{Verification status}

\begin{center}
\begin{tabular}{@{}lp{9.5cm}@{}}
\toprule
Aspect & Status \\
\midrule
Registry, dummy path & exercised by the test suite \\
State loader & shape-level assertion only \\
Environment loader & shape-level assertion only \\
End-to-end run & exercised with the dummy model: file exists, sidecar
  exists and parses, raster opens with a CRS \\
Real checkpoint load & \textbf{not} covered \\
Inference correctness & \textbf{not} covered \\
National and continental aggregators & \textbf{not} covered \\
Band names, sidecar contract label & \textbf{not} asserted \\
Degradation to zeros & \textbf{not} asserted \\
\bottomrule
\end{tabular}
\end{center}

The coverage described here is what the repository's test suite exercises
around this stage; the assertions it makes are about existence and shape
rather than about values. No test in the repository loads a real checkpoint
and checks the inference; no test exercises the national or continental
aggregators; no test asserts the output's band names or the sidecar's
contract label; and no test asserts that a missing input produces a warning
or a distinguishable output. The stage's end-to-end path is therefore
verified only in the sense that a raster and a sidecar appear.

\subsection{What is verified and what is not, in this document}

Verified against the code by direct reading: the band names and their
meanings; the contract versions; the order of operations; the fallback
chains; the aggregation rules; the output profile; the sidecar's key list;
the behaviour of the scenario parameter; the absence of version gating; the
absence of any classification, calibration or uncertainty treatment; and the
absence of a manifest consumer.

Verified by inspection of the tree: no trained checkpoint and no
\code{model.yaml} exists anywhere in the repository --- the runs tree holds
no \code{*.pt} and no \code{model.yaml} --- so there is no checkpoint under
the current state contract, and none under a retired one either.

Not established, and stated as such in the body: whether the national block
mean or the description string is the intended national semantics; whether
the transform/dimension discrepancy at the coarser scales is a defect or the
intended contract; which environment channel order is intended; and whether
the transmission raster and the vector snapshot are ever intended to come
from the same day.

\subsection{Statements about the retired specification}

The retired specification for this stage disagrees with the implementation in
several places, and this document follows the implementation. Where the two
differ on the number of input channels, on the number of state channels, on
the day clamp, on the band description, on the source of the NoData value and
on the existence of an AOI catalogue symbol, the code is authoritative. The
specification's own description of the stage's consumers --- maps,
dashboards, intervention planning --- describes consumers that do not exist
in the tree.

\end{document}
